% ...........................................................................
% 04c-reencoding.tex -- beats 5-6. The residual re-encoding.
% Legacy anchors sec:methods-vardecomp and sec:methods-holdout are re-\label-ed
% here so v1 cross-references resolve.
% ...........................................................................

\beatsection{R}{5}{The standard target hides the drug}{sec:methods-residual}

\begin{question}[5]
  Every target compared in the last section ranks genes by how much of them a cell
  holds, and two drugs on the same cells hold almost the same genes in almost the
  same order, so the loss can barely tell those targets apart. The apparatus chapter
  defines a target that ranks something else instead. Does it in fact expose the
  drug, and can it be estimated without the evaluation leaking into the fit? The
  instrument counts how many top target positions carry a different gene for two
  drugs in one context. 
  \end{question}

  \noindent
  \Cref{sec:res-epsilon} ended on a distinction it did not resolve: producing no
  structure and producing structure that does not track the drug are different
  failures. This section takes up the second, and it opens by conceding the first.
  Three target constructions were tried there and none separated, so a fourth
  arrives owing an explanation of why it is not simply a fourth point on the same
  ladder.
  
  It is not, because that ladder held one thing fixed. Every target in
  \cref{tab:epsilon} ranks genes by how much of them the cell holds and varies only
  how many cells are averaged before the ranking is taken. Averaging changes how
  noisy a ranking is; it does not change what the ranking is \emph{of}. Two drugs
  applied to the same cells leave the same abundant genes at the top in nearly the
  same order, so every point on that ladder inherits nearly the same target string
  however cleanly it is estimated --- which is what \cref{sec:res-epsilon} meant in
  reserving judgement on constructions outside its family.
  \Cref{sec:apparatus-vocab} defines one. What follows is the measurement that
  motivated it and the analysis of whether it can be estimated at all.
  
  \emph{Encoding} here means how the target string is written, not what the model's
  internal states carry (\cref{sec:res-used}). The loss weights every position
  equally, so the measurement is a count over the $200$ target positions
  (\cref{tab:divergence}).
  
  % ...........................................................................
% ...........................................................................
\begin{table}[htbp]
  % MARGINALIA: full width; the long caption sets above at the text width,
  % notes below in the sans. One rule, under the head.
  \caption{\captiontitle{$83\%$ of the standard target's tokens are identical between any
  two drugs in the same context; under the residual, $40\%$ are --- and the
  residual's exposure does not track how similar the two drugs' responses are.}
  Rows are quoted at the plate scope this build uses (\cref{sec:apparatus-vocab});
  at cell-line scope the residual row reads $118.2$ overall and $0.7432$, and the
  other two rows are unchanged. The three stratum columns bin all $\num{212528}$
  exhaustive within-group drug pairs on the cosine between the two residuals ---
  near above $+0.1$ ($n = \num{23024}$), orthogonal within $\pm 0.1$
  ($n = \num{152421}$), opposite below $-0.1$ ($n = \num{37083}$) --- on the same
  similarity axis and the same cut points the later swap evaluation uses. The full
  profile and the shift climb monotonically across those strata and correlate
  negatively with residual cosine; the residual stays high in all three and does
  neither. That weakens the reading that the later swap-distance pattern is
  inherited trivially from target divergence rather than being an independent test
  of it. The counts are where the three encodings differ, and they do not on their
  own identify tokenisation as the binding constraint.}
  \label{tab:divergence}
  \begin{fullwidth}
  \begin{threeparttable}
  \begin{tabular}{@{}l r@{\hspace{0.4em}}l r r r r r@{}}
  & \multicolumn{5}{l}{\thead{Top-$200$ tokens differing between drugs}}
  & & \\
  \addlinespace[2pt]
  \thead{Representation} & \thead{all pairs} & & \thead{near} & \thead{orth.}
    & \thead{opposite} & \thead{$\rho$}
    & \thead{Retrieval reference} \\
  \midrule
  full profile (the standard target)
    & $34.6$ & ($17\%$)  & $32.5$ & $34.6$ & $35.8$ & $-0.149$ & $0.742$ \\
  shift (treated $-$ control)
    & $111.4$ & ($56\%$) & $107.4$ & $111.3$ & $114.5$ & $-0.282$ & $0.742$ \\
  residual (shift $-$ generic)
    & $120.3$ & ($60\%$) & $118.4$ & $121.0$ & $118.7$ & \key{$+0.019$} & $0.7467$ \\
  \end{tabular}
  \begin{tablenotes}[flushleft]
  \item The count columns and the retrieval column hold different quantities, a
  tally out of $200$ positions and a retrieval score, so rows compare within a
  column and never columns with each other. $\rho$ is the Spearman correlation
  between a pair's residual cosine and its differing-token count, taken over all
  pairs and not within a stratum.
  \item The retrieval column pits disjoint cell subsets of one drug, in one cell
  line, on one plate, against each other. It is a precision reference, not a
  biological replicate ceiling, and it is near-constant across the three rows: the
  rewrite exposes the drug without buying that exposure from a coarser measurement.
  \item The stratification is descriptive. Pairs reuse drugs and experimental
  groups, so no pair-level inferential claim is made and no interval is drawn.
  \end{tablenotes}
  \end{threeparttable}
  \end{fullwidth}
  \end{table}
  \paragraph{Which subtraction does the work} Removing the control moves the
  between-drug token difference from $34.6$ to $111.4$; removing the generic adds a
  further $8.9$ at the plate scope this build uses, and $6.8$ at cell-line scope. So
  the control subtraction alone buys $90\%$ of the recoverable difference, $92\%$
  cell-line-scoped. That is the first thing the table says and the least comfortable
  one for the construction it motivates: almost all of the exposure comes from the
  cheaper of the two subtractions. The measurement is an overlap of gene sets on one
  atlas at $K = 200$ ($\num{6602}$ conditions, \texttt{target\_divergence.json}), so
  no margin is claimed for other atlases or other $K$. It squares with
  \cref{sec:res-metrics}, where what saturates $\DEdr$ is reversion toward a central
  point, \texttt{revert\_center} scoring $0.9999$ without fitting anything.
  
  Read as a verdict on the generic, though, that $90\%$ proves less than it looks
  like, and the reason is in what the count counts. It tallies how many of the $200$
  positions carry a different gene between two drugs, so a subtraction is credited
  only where it reorders the two drugs' rankings differently \emph{from each other}.
  Removing a component the two drugs share does that weakly by construction, however
  large the component is. As we show further in later sections however, by magnitude the split runs the other way: the shift's
  squared norm decomposes into a drug-specific residual of $0.621$, a generic
  programme of $0.390$ and a cross term of $-0.011$. The generic is two fifths of
  what the control subtraction leaves behind, and it is very nearly orthogonal to
  the part that identifies the drug.
  
  Thus the two subtractions are not competing for one credit, then. The control
  subtraction is what exposes positions to the drug at all, and the generic
  subtraction decides what those exposed positions are about --- a target that
  skipped it would rank each drug's genes against a baseline still carrying two
  fifths of the programme every drug triggers. Neither is dispensable, and
  \cref{tab:divergence} is simply the wrong instrument for asking whether the second
  one is.
  
  One threshold choice belongs beside these counts. The inclusion minima of
  \cref{sec:apparatus-vocab}, at least $40$ treated and $20$ plate-matched control
  cells, sit where the identifiability curve of \cref{sec:res-blind} flattens: the
  within-well split-half precision reference reads $0.61$ at $10$ cells, $0.76$ at
  $40$ and $0.85$ at $120$.
  
  \paragraph{The count stops short of the tokenisation claim} The natural inference,
  that only a small fraction of the gradient carries drug identity, is a hypothesis
  about the optimisation; what is measured is the overlap of target gene sets, not a
  token-level loss or a gradient attribution. Gene positions are not loss-token
  positions under byte-pair encoding, and the residual target alters aggregation,
  control referencing, generic subtraction, sign coding and sequence length at once.
  Separating those five factors needs a matched-budget factorial this thesis does not
  run.
  
  
  % ...........................................................................
  \paragraph{What the plate-local generic costs, measured}
  \Cref{sec:apparatus-vocab} states the limit as part of the definition: what
  survives both subtractions is not what makes this drug different from all drugs,
  but what makes it different from the other training drugs sharing its plate and
  cell line. Stated there it is a definitional caveat. Here it is a quantity, and
  mechanism is where it bites. \Cref{sec:res-scope} finds mechanism-related drugs
  co-plated $0.404$ of the time against $0.362$ for a random draw, so for a co-plated
  mechanism class part of the class's own shared programme sits inside the generic
  and is subtracted out of every member's residual by construction.
  
  That bias was tested rather than assumed, and the mechanism gate of
  \cref{sec:res-scope} is where it can be, because that gate is scored in this
  frame.

  \begin{scopelimit}[carried into \cref{sec:res-scope} and \cref{sec:limitations}]
  Two properties of this frame travel with every number measured in it. First, the
  generic is a plate-local mean, so where a mechanism class is co-plated the
  mechanism channel is later graded in a frame that has already removed some of what
  it is asked to find, which makes an inconclusive verdict on mechanism the
  conservative reading and not a generous one. Second, the plate-scoped generic
  removes plate structure from the truths and, in expectation, from the predictions,
  but that is an expectation and not an identity. An individual residual may retain
  some plate structure, and the size of what remains was never measured. The absolute
  $\NIR$ values of this frame are read within it and never set beside the
  within-plate tables elsewhere in the chapter.
  \end{scopelimit}
  
  % ...........................................................................
\paragraph{Estimating the generic is where leakage would enter} The residual is a
subtraction, and the quantity subtracted is not handed down by the design: it is
estimated from the same data the model is trained and scored on. That makes the
generic the one part of the construction capable of going wrong in a way the
target string never shows. This is particularly concerning if we wish to perform preprocessing on conditions and admit only conditions that carry signal and not noise. A leaky generic still produces well-formed sentences of
plausible genes; it simply produces them against a baseline that has already seen
what the model is about to be tested on. \Cref{sec:apparatus-vocab} states two
properties of $D_c$ as fundamental, leave-one-drug-out and split-before-fit,
without arguing for either. They are argued here, because between them they settle
the two questions an estimated subtrahend raises: what the mean is taken over, and
which conditions are allowed into it.

\paragraph{The mean is over drugs, and a condition's own drug leaves it entirely}
The generic is built in two steps rather than one. Within a group, each drug's
conditions are first averaged into a single per-drug shift; the generic is then the
mean of those per-drug shifts. Both steps carry weight. Averaging within a drug
first is what makes the result a mean over drugs at all: otherwise a drug measured
at five doses in the group would contribute five terms and count five times as
heavily as a single-dose drug, and the generic would be a mean over conditions
wearing the name of a mean over drugs.

\emph{Leave-one-out} then names a removal that could be drawn at several levels,
and the level is what the rule is about. What is removed is the drug. When the
generic for a condition is formed, that condition's drug is struck from the group's
roster altogether --- every dose of it and every condition of it, not merely the one
condition whose residual is being computed --- and the mean is taken over the drugs
that remain. The narrower alternative, dropping only the condition itself, sounds
equivalent and is not. A drug run at five doses would keep four of its own
conditions inside the mean subtracted from its fifth, so part of what that drug does
would be subtracted from itself, and the residual would answer a different question:
what makes this dose unlike this drug's other doses, rather than what makes this
drug unlike the other drugs. Removing at drug level is what guarantees that the
quantity subtracted from a condition carries no information about that condition's
drug.

The rule also keeps the two splits commensurable, which is the reason less often
noticed. The generic is fitted on training conditions only, so a training
condition's drug sits in the fit inventory and, without the rule, would appear in
the mean subtracted from itself. A held-out condition's drug contributes nothing to
its group's generic in any case, because the split holds out every condition of that
(drug, cell line) pair. The two splits would therefore be scored against
differently defined targets, and a comparison between them would be confounded by
the definition rather than by generalisation. Leave-one-drug-out removes the
asymmetry by making both sides alike: neither has its own drug in its own generic.

\paragraph{Which conditions may enter the generic, and what makes the question
hard} The second question has an answer that raises no difficulty whatever: every
condition the split assigns to training. Fitting the generic on that set is one
pass over fixed data, nothing in it refers to any quantity computed later, and
there is no circularity to speak of. The difficulty is introduced deliberately,
by us, and it is introduced for a reason worth keeping.

Not every condition carries a residual worth training on. At this cell budget the
two subtractions can leave a vector that is almost entirely sampling noise, and
such a condition still emits a full-length target string --- ranked, closed with
\ENDCELL{}, indistinguishable in form from any other --- which the loss weights
exactly as heavily as a condition whose residual is real. Training on it teaches
the model to reproduce noise in the voice of a drug effect. The scale of that is
not marginal: the line finally adopted retains half the training conditions, which
is another way of saying that the median condition's two halves agree about as
well as two unrelated conditions' halves do. A filter is worth having, and the
natural test for one is self-consistency.

That test is the half-split, and its content is precise. A \emph{condition} ---
one drug, at one dose, on one cell line, on one plate, the unit of
\cref{sec:methods-data} --- has its treated cells shuffled and cut in two, and the
plate's control cells are cut in two alongside them; the same halving is applied
to every other condition in the same (cell line, plate) group. Each half yields a
shift, its treated mean minus its control mean, computed inside that half alone.
Each half also yields a generic, the leave-one-drug-out mean just described taken
over the other training drugs' shifts on the matching half. Subtracting the one
from the other leaves two residuals that share no cell and no estimate, and the
condition is kept only if the cosine between them clears a threshold. Two
independent measurements of the same condition must agree before it is allowed to
teach anything.

Adding that test is what closes the loop. It decides membership of the training
set, and it decides it from the residual --- which is a shift minus the generic,
and the generic is a mean over the other \emph{training} drugs. Membership is
therefore settled by a quantity that is itself a function of membership. Note
where the problem is not: nothing in the test is recursive, since a shift needs
only cell means and a generic needs only other drugs' shifts, so for any fixed
list of training drugs every quantity computes in a single pass and a condition's
residual never requires its own residual. The circularity sits one level up and it
is about \emph{membership} alone. Written as edges: the reliability line selects
the training set, the training set fits the generic, the generic is subtracted to
give the residual, and the residual is what the reliability line is applied to
(\cref{fig:circularity}, left). Every edge is defensible on its own, the four
close, and nothing in the definitions says where to start.

What that would cost is not abstract. Suppose the generic were refitted on the
conditions the filter retains, which is the natural reading of ``the mean over the
drugs we are actually training on''. Dropping one condition would then change the
generic of every other condition in its group, hence their residuals, hence which
of \emph{them} pass --- so whether a condition survived would depend on when it
was evaluated, and the procedure would have to be iterated to a fixed point that
nothing guarantees to exist or to be unique.

So the design problem is not whether to filter but where to let the filter act. It
must be allowed to decide what the model trains on, which is the whole point of
it, and it must not be allowed to decide what the generic is fitted on, which is
what would close the circle. Those two can be separated, because they are separate
sets: the fit inventory is fixed by the metadata split, and the filter operates
downstream of it on conditions whose residuals have already been computed. The
generic is fitted once, on every condition the split assigns to training, before a
single reliability test has run; the filter is one terminal pass against that
fixed estimate; and dropping a condition moves no other condition's residual. The
retained set is a function of the data and not of the order it was walked in.

The hazard is set out at this length because an earlier build of ours was in it,
estimating generic and filter over every condition and assigning the holdout
afterwards. That adds a fifth edge to the figure, held-out conditions entering the
fit, and makes each held-out target partly a function of the other held-out
conditions. The failure is quiet: it survives every check that inspects the
targets themselves, since the targets it produces look exactly like the ones it
should have produced. Any construction that estimates what it subtracts, and then
filters on what is left, is exposed to some version of it.

\paragraph{The loop is cut by ordering, and the ordering is enforced} The build runs
as \emph{inventory}, \emph{split}, \emph{fit}, \emph{transform}, and each stage may
read only what the stage before it is allowed to expose (\cref{fig:circularity},
right). Inventory enumerates conditions from metadata and cell counts alone --- the
(drug, cell line, plate, dose) tuples clearing $40$ treated and $20$ plate-matched
control cells --- and aggregates no expression at all. Split then assigns train and
holdout from that metadata, keyed on the (drug, cell line) pair so that every dose
and every plate of a held-out pair leaves with it, and it does so before a single
pseudobulk exists. That removes the first edge: nothing downstream can have defined
the split, because at the moment the split is written there is nothing downstream to
define it with.

Fit estimates the generic, and is handed the training keys explicitly rather than
inferring them. Its inventory is every condition the split assigned to training,
taken whole and before any reliability test has run, which removes the second edge
--- admission cannot select the fit inventory, because the fit inventory is fixed
before admission is decided. Within that inventory the estimator is the one argued
above, drug-weighted and leave-one-drug-out, scoped to the (cell line, plate) group
and failing closed where a group holds fewer than three other training drugs, in
which case the condition is dropped and counted rather than quietly falling back to
a coarser scope. Transform then projects every condition, held-out ones included,
through that fitted generic. Held-out conditions are scored by it and never enter
it.

The training set is made last, and in a single pass. Each training condition's two
half-residuals are formed against the fixed generic, their cosine is taken, and the
condition is kept if it clears the line, $2{,}523$ of the $\num{4999}$ training
conditions or $50.5\%$. Because the generic is never refitted on the survivors, a
drop moves no other condition's residual and the retained set does not depend on the
order the pass walked in. Held-out conditions are transformed but deliberately not
filtered. Each retained condition then emits one training example per plate-matched
control cell, the prompt carrying that control cell together with the drug, its dose
and its mechanism annotation, and the response being that condition's residual
string.

The staging is checkable rather than merely asserted. A poison test corrupts every
held-out expression vector with large noise, rebuilds, and requires both the
training \texttt{JSONL} and a hash of every fitted quantity to be
unchanged.\gloss{poison test}{Corrupt the held-out side; the training targets and
the fit digest must not move. Corrupt a training condition; they must.} If one
held-out cell reaches the fit, one of the two moves, and restoring the earlier
build's one-line defect does make it fail on the digest, so the check is not
vacuous. Poisoning a training condition instead must change both, which is the
negative control that stops it passing for the wrong reason. One quantity the
staging does not settle is the threshold the filter applies, and that is the last
piece.

% ===========================================================================
% circularity.tex -- the residual definition's loop, and the ordering that cuts
% it. \input by Sections/04c-reencoding.tex at sec:methods-residual. Compiles
% against thesis_v2/preamble.tex and nothing else: no external data, no
% \includegraphics, no package this document does not already load
% (tikz + arrows.meta, both in preamble.tex section 9).
%
% THIS IS A PROVENANCE DIAGRAM, NOT A RESULT. Deliberately absent: every NIR
% value, every arm, every split-level number. Two counts appear, once, and both
% are already in the prose of this same section. Sibling circularity.json lists
% every number drawn.
%
% THE LEFT PANEL CARRIES NO BUILD LABEL, DELIBERATELY. It drew "THE EARLIER
% BUILD -- NOT THE SHIPPED ONE" until this revision, which spent half a
% full-width figure on a pipeline the thesis did not run. The historical fact
% -- that an earlier build did close this loop -- is a disclosure, and a
% disclosure belongs in prose where it can be qualified: it is stated at
% Sections/04c-reencoding.tex:148-151 ("That is a loop [...], and an earlier
% build closed it, with generic and filter estimated over every condition and
% the holdout assigned afterwards"). That sentence stands on its own and was
% NOT deleted when the label came off the panel. What the left panel draws now
% is the structural hazard itself, in the present tense and attached to no
% build: the cycle the four definitions form, and the edge that fitting before
% splitting adds to it. The two panels therefore read hazard and remedy, which
% is the only reason the staged ordering is load-bearing rather than fastidious.
%
%   2,523 of 4,999 training conditions retained (50.5%)   04c-reencoding.tex:211-212
%   the resolved reliability line, +0.1086                 04c-reencoding.tex:211
%
% Cross-checked, all four agreeing, before drawing:
%   Sections/03-Apparatus.tex:206     "2,523 of 4,999 train (50.5\%)"
%   Sections/04c-reencoding.tex:212   "$2{,}523$ of $\num{4999}$ ... ($50.5\%$)"
%   Sections/04d-baselines.tex:337    "only $\num{2523}$ are training"
%   Sections/05-Limitations.tex:33    "$\num{2523}$ of $\num{4999}$ ... $\cos > 0.109$"
% +0.109 is the measured null's 95th percentile and +0.1086 is what it resolves
% to; the prose prints both, and the node carries the resolved value because
% that is the one the builder applies.
%
% LAYOUT NOTES, all forced by a build:
%   * fullwidth = 174mm = 493.23pt (textwidth 142mm + marginparsep 4mm +
%     marginparwidth 28mm, preamble.tex:88-90 and :525). The bounding box is
%     PINNED to 17.35cm = 491.81pt, 1.4pt inside the measure. Unpinned, the
%     picture measured 3.10pt OVER and the figure body reported an overfull
%     hbox: nodes with anchor=north west hang their 3.33pt inner sep left of
%     the anchor at x=0.20, and the strip rules end at x=17.35. Pinning clips
%     no ink -- no drawn object lies outside [0, 17.35] x [-0.95, 7.95].
%   * NO node uses text width. Every line break is written by hand. The first
%     build did use text width and TeX hyphenated inside the boxes
%     ("the relia-/bility line"), grew four of them past their minimum heights
%     and overlapped the rows below. Hand-broken lines make every box width and
%     height deterministic, which is what lets the coordinates below be trusted.
%   * the four objects are ONE node set drawn twice, same style, same wording,
%     same 2.80cm minimum width, so the right panel reads as a re-laying of the
%     left and not as a second diagram. Only the arrangement changes.
%   * the fatal edge is drawn and STRUCK, not deleted. A deleted edge leaves
%     nothing for the caption to point at, and the passage is about which single
%     edge the ordering removes.
%   * the two streams leaving SPLIT must not cross. They cannot both turn upward
%     in the order they leave: the stream that leaves HIGHER (held-out, 2.2mm
%     above the east anchor) turns up FIRST (x=10.30) and the stream that leaves
%     LOWER (training keys, at the east anchor) turns up LATER (x=10.85). Any
%     other pairing crosses, and a crossing here reads as contact between them.
%   * both streams use -| rather than an explicit corner coordinate, so the
%     corner inherits the node anchor's own y. Written as a literal corner, the
%     first segment came out as a visible diagonal whenever a node's height
%     differed by a hair from the value assumed at design time.
%   * the held-out stream is routed over the top of the fit column, and its one
%     arrowhead points right, into TRANSFORM. No arrowhead anywhere in the right
%     panel points back left or up into the fit; that absence is the claim.
%   * NO edge label is masked onto its arrow. The first build masked all of
%     them and the labels ate the edges: the C->D gap in the left panel is
%     1.20cm and "subtracted to give" is 1.15cm wide, so the bottom edge of the
%     cycle disappeared and the loop stopped reading as a loop. Every label now
%     sits beside its edge. The tightest is "training keys", hung under the
%     corner at x = 10.43 to 11.27, between the held-out riser at 10.30 and the
%     fit column's left edge at 11.35.
% ===========================================================================
\begin{figure}[htbp]
\begin{fullwidth}
\centering
\begin{tikzpicture}[x=1cm, y=1cm]
  \useasboundingbox (0,-0.95) rectangle (17.35,7.95);
  \tikzset{
    % the shared node set. Identical style in both panels.
    obj/.style={draw=rulegrey, line width=0.4pt, fill=white, inner sep=3pt,
                align=center, font=\scriptsize, text=ink,
                minimum width=2.80cm, minimum height=0.68cm},
    edg/.style={-{Stealth[length=3.6pt,width=2.8pt]}, draw=ink,
                line width=0.5pt},
    lab/.style={font=\scriptsize, text=quietink, fill=white, inner sep=1pt,
                align=center},
    side/.style={font=\scriptsize, text=quietink, align=flush right,
                 anchor=east, text width=3.20cm},
    quiet/.style={font=\scriptsize\itshape, text=quietink, align=center},
    tag/.style={font=\sffamily\scriptsize\bfseries, text=accent},
    sub/.style={font=\footnotesize\itshape, text=quietink}}
  \def\T#1{{\sffamily\scriptsize\bfseries\color{accent}#1}\\[1.0pt]}
  \def\Q#1{\\[1.0pt]{\color{quietink}#1}}

  % ======================= PANEL HEADS =====================================
  % Hazard and remedy, one head each, same shape: a tag naming the structure
  % and a sub saying what happens to the four objects. Neither names a build.
  \node[tag, anchor=north west] at (0.20,7.90) {THE CIRCULARITY};
  \node[sub, anchor=north west] at (0.20,7.55)
    {the four objects define one another};
  \node[tag, anchor=north west] at (7.50,7.90) {THE ORDERING THAT CUTS IT};
  \node[sub, anchor=north west] at (7.50,7.55)
    {the same four objects, staged by execution order};

  \draw[rulegrey, line width=0.4pt] (7.20,1.15) -- (7.20,6.60);

  % ======================= LEFT: THE CLOSED LOOP ===========================
  \node[obj, anchor=north] (A) at (1.60,6.10) {the reliability line};
  \node[obj, anchor=north] (B) at (5.60,6.10) {the training set};
  \node[obj, anchor=north] (C) at (5.60,4.45) {the generic};
  \node[obj, anchor=north] (D) at (1.60,4.45) {the residual};

  % NONE of the four loop labels is masked onto its arrow. Masked, they ate the
  % edges: the C->D gap is 1.20cm and "subtracted to give" is 1.15cm wide, so
  % the bottom edge of the cycle vanished entirely and the loop stopped reading
  % as a loop. Each label now sits beside its own edge -- the two horizontals
  % labelled outside the cycle, the two verticals inside it -- and all four
  % edges draw unbroken.
  \draw[edg] (A.east) -- (B.west);
  \draw[edg] (B.south) -- (C.north);
  \draw[edg] (C.west) -- (D.east);
  \draw[edg] (D.north) -- (A.south);
  \node[lab, anchor=south] at (3.60,5.86) {selects};
  \node[lab, anchor=east]  at (5.45,4.94) {fits};
  % clears the box row's bottom edge at y=3.77: at y=3.97 its first line sat
  % level with the two boxes' lower corners, 0.02cm from each, and read as a
  % collision even though it was not one.
  \node[lab, anchor=north] at (3.60,3.88) {subtracted\\to give};
  \node[lab, anchor=west]  at (1.75,4.94) {must reproduce\\across a half-split};

  % --- the fatal edge: drawn once, struck once, and labelled --------------
  % Present tense, and attached to an ordering rather than to a build: this is
  % what fitting before splitting does to the cycle, not what one run once did.
  % The quiet italic under H answers "why would they enter the fit?" inside the
  % panel, and mirrors the right panel's one quiet italic under SPLIT, so each
  % half carries exactly one such annotation.
  \node[obj, anchor=north] (H) at (1.60,2.60) {held-out conditions};
  \draw[edg] (H.east) -| (C.south);
  \node[lab, anchor=east] at (5.45,2.95) {enter the fit};
  % Its y is the strike caption's y, not a value of its own: both are two-line
  % blocks of the same size, so sharing anchor y = 1.74 gives them a common top
  % AND a common bottom (0.944) and the bottom-left reads as one row rather than
  % two blocks at two heights. Widest line measures 54.42pt = 1.91cm, so with
  % inner sep the node spans x 0.53-2.67, clear of the strike caption's left
  % edge at x = 3.10 by 0.43cm. Top sits 0.18cm under H, matching the 0.20cm
  % the right panel leaves under SPLIT.
  \node[quiet, anchor=north] at (1.60,1.74)
    {when the fit\\precedes the split};
  % the strike. One stroke. Nothing else in this figure is accent-coloured
  % except the structural tags.
  \draw[accent, line width=1.1pt] (4.25,1.94) -- (4.75,2.58);
  \node[font=\sffamily\scriptsize, text=accent, anchor=north west,
        align=left] at (3.10,1.74)
    {struck: the one edge\\the ordering removes};

  % ======================= RIGHT: THE STAGED PIPELINE ======================
  \fill[panelbg] (7.45,3.55) rectangle (10.15,6.30);
  \node[obj, anchor=north, minimum width=2.55cm] (INV) at (8.80,6.10)
    {\T{INVENTORY}wells, drugs, doses};
  \node[obj, anchor=north, minimum width=2.55cm] (SPL) at (8.80,5.10)
    {\T{SPLIT}train \textbar{} held out};
  \draw[edg] (INV.south) -- (SPL.north);
  \node[quiet, anchor=north] at (8.80,4.22)
    {metadata only:\\no expression data};

  % FIT, and the training side that descends from it
  \node[obj, anchor=north] (G) at (12.75,6.10) {\T{FIT}the generic};
  \node[obj, anchor=north] (R) at (12.75,4.60) {the residual};
  \node[obj, anchor=north] (L) at (12.75,3.30)
    {the reliability line\Q{$+0.1086$}};
  \node[obj, anchor=north] (S) at (12.75,1.90) {the training set};
  \draw[edg] (G.south) -- (R.north);
  \draw[edg] (R.south) -- (L.north);
  \draw[edg] (L.south) -- (S.north);

  \node[side] at (11.25,2.80) {$2{,}523$ of $4{,}999$ training conditions
    retained ($50.5\%$)};
  \node[side] at (11.25,1.56) {the residual targets the model is trained on};

  % training keys, handed to the fit explicitly
  % -| , not an explicit corner coordinate: the corner must inherit SPL.east's
  % own y, or the first segment comes out as a visible diagonal when the node's
  % height differs by a hair from the value assumed here.
  % This is the one edge label NOT masked onto its arrow. The riser is 1.0cm
  % long and a masked label ate 0.6cm of it, leaving the arrow looking like two
  % disconnected stubs. It hangs under the corner instead, in the only clear
  % space in the panel: x 10.43-11.27, below the horizontal and above the
  % retention annotation, which tops out at y = 3.25.
  \draw[edg] (SPL.east) -| (10.85,5.675) -- (G.west);
  \node[font=\scriptsize, text=quietink, align=center, anchor=north]
    at (10.85,4.50) {training\\keys};

  % TRANSFORM, fed by the fitted generic and by a stream that bypasses the fit
  \node[obj, anchor=north] (TR) at (15.95,6.10)
    {\T{TRANSFORM}held-out conditions\\projected through\\the fitted generic};
  \draw[edg] (G.east) -- (G.east -| TR.west);
  \draw[edg] ([yshift=2.2mm]SPL.east) -| (10.30,6.45) -- (15.95,6.45)
    -- (TR.north);
  \node[font=\scriptsize, text=quietink, anchor=south] at (13.10,6.55)
    {held-out keys --- they never enter the fit};

  \node[obj, anchor=north] (HT) at (15.95,4.30)
    {held-out residual\\targets};
  \draw[edg] (TR.south) -- (HT.north);
  \node[quiet, anchor=north] at (15.95,3.25)
    {deliberately unfiltered};
  \node[quiet, anchor=north] at (15.95,2.80)
    {nothing here\\re-enters the fit};

  % ======================= THE POISON TEST STRIP ===========================
  \draw[rulegrey, line width=0.4pt] (0.20,0.85) -- (17.35,0.85);
  \node[tag, anchor=north west, align=left] at (0.20,0.46)
    {THE POISON\\TEST};
  \node[font=\sffamily\scriptsize\bfseries, text=ink, anchor=north west]
    at (3.30,0.55) {CORRUPTION APPLIED};
  \node[font=\sffamily\scriptsize\bfseries, text=ink, anchor=north west]
    at (9.60,0.55) {REQUIRED OUTCOME};
  \draw[rulegrey, line width=0.4pt] (3.30,0.20) -- (17.35,0.20);
  \node[font=\sffamily\footnotesize, text=ink, anchor=north west]
    at (3.30,0.08) {corrupt every held-out expression vector};
  \node[font=\sffamily\footnotesize, text=ink, anchor=north west]
    at (9.60,0.08) {the training targets are byte-identical};
  \node[font=\sffamily\footnotesize, text=ink, anchor=north west]
    at (3.30,-0.40) {corrupt one training condition};
  \node[font=\sffamily\footnotesize, text=ink, anchor=north west]
    at (9.60,-0.40) {the training targets must move};
  \draw[rulegrey, line width=0.4pt] (3.30,-0.85) -- (17.35,-0.85);
\end{tikzpicture}
\caption[The residual's circular definition, cut by
ordering]{\captiontitle{The definition of the residual could have closed on itself,
and what stops it is execution order, not an argument.} \emph{Left, the
circularity.} A condition enters training only if
its residual reproduces across a half-split; the residual exists only once the
generic is subtracted; the generic is fitted on the training set; and the
training set is what the reliability line selects. Estimate the generic and the
filter over every condition and assign the holdout afterwards, and held-out
conditions enter the fit --- the struck edge --- so each held-out
target is defined partly by other held-out conditions. \emph{Right, the
ordering that cuts it,} the build this thesis reports. The same four objects,
staged. Inventory and split read
metadata only, before any pseudobulk exists, so nothing downstream of the split
can have defined it; the fit is handed the training keys explicitly; and
transform projects held-out conditions through the fitted generic in one
direction, with no path back into it. The generic is the mean over training
drugs other than this one, within plate. The two counts are the
\emph{training-set} retention at the resolved reliability line $+0.1086$, the
$95$th percentile of a null pairing half $A$ of one condition with half $B$ of
another; the held-out sets are deliberately unfiltered, since dropping held-out
conditions for failing their own reproducibility bar would select the
evaluation set on its own outcome. \emph{Beneath,} the two assertions that make
the ordering checkable in both directions rather than merely asserted. This is
a provenance diagram: no $\NIR$ value, arm or split-level result appears in
it.}
\label{fig:circularity}
\end{fullwidth}
\end{figure}

  \paragraph{How the reliability line was chosen} The threshold decides more of the
training set than any other single choice, so it is measured rather than picked.
The builder constructs a matched null by pairing half $A$ of one condition with
half $B$ of a \emph{different} condition --- same encoding, same dimensionality,
same noise level, no shared biology --- reports the whole retention curve against
it, and places the line at that null's $95$th percentile, $+0.109$. What that implements is a criterion
rather than a setting. A condition is kept when its two halves agree more than two
unrelated conditions' halves do, the line being the level of agreement only one
unrelated pair in twenty reaches. It retains $2{,}523$ of the $\num{4999}$
training conditions, $50.5\%$.

The filter is applied to the training side alone. Dropping held-out conditions for
failing their own reproducibility bar would select the evaluation set on its own
outcome, and the two errors are not symmetric: noise in a retained training
condition costs learning efficiency, whereas noise admitted to the evaluation
would cost validity.

What clearing the line certifies is sampling precision and not biology. Both
halves are drawn from the same well, so a condition that clears it has a residual
estimated stably at this cell budget, which is not the claim that the effect would
reappear in a fresh experiment. $\num{286}$ of the $\num{287}$ distinct (drug,
dose) treatments in this cache occur in exactly one well, so a biological
replicate is not available to ask for.
% ===========================================================================
% fig-frame.tex -- v2. The residual frame, the encoding fork, and the SCORED
% object. \input by Sections/04c-reencoding.tex. Compiles against
% thesis_v2/preamble.tex and nothing else; no external data, no
% \includegraphics. FULLWIDTH, and MEASURED rather than assumed: the
% tikzpicture's own bounding box is 491.31pt wide against a fullwidth measure
% of 495.08pt (\textwidth 404.03pt + \marginparwidth + \marginparsep). The
% measurement is made by the picture itself, guarded by \fframedebug;
% figs/_harness_frame.tex switches it on, and it is inert in main.tex.
%
% RULE INHERITED FROM preamble.tex section 8: a page carrying a fullwidth float
% carries NO \marginpar. Every note in this figure is therefore drawn INSIDE
% the tikzpicture. Nothing here calls \gloss, \seealso or any margin wrapper.
%
% EVERY NUMBER IN THIS FIGURE, WITH ITS SOURCE. Nothing is invented.
%   34.6 / 200  (17%)  full profile, cell-line scope
%   118.2 / 200 (59%)  residual,     cell-line scope
%   120.3              residual,     plate scope            -- caption only
%   0.742 -> 0.743     within-well split-half retrieval reference, full ->
%                      residual (stored 0.741986901358855 and 0.7432424792)
%     RESULTS_cluster/target_divergence.json, scopes.cellline.{full,residual}
%     .divergence_mean and .ceiling_retrieval, and scopes.plate.residual;
%     and tab:divergence, Sections/04c-reencoding.tex lines 42-44.
%   100 genes per block, [DOWN], [END_CELL], weight 1/log2(rank+1), cosine,
%   NIR over other drugs in the same cell line
%     Sections/04c-reencoding.tex, "How the residual becomes a string, and how
%     it is scored"; and endcell/analysis/residual_eval.py,
%     signed_rank_from_sentence / signed_rank_from_vector (k = 100,
%     v[gi] = sgn / np.log2(r + 1.0)).
%   946 panel genes -- Sections/04c-reencoding.tex, "a 946-dimensional profile".
%   the stem weights: 1/log2(r+1) for r = 1..8 =
%     1.0000, 0.6309, 0.5000, 0.4307, 0.3869, 0.3562, 0.3333, 0.3155.
%     Every stem in stage 3 is one of these eight numbers; the +1 / -1 ticks
%     are there so a reader can check them.
%   "fewer than three other training drugs yields no generic and the condition
%     is dropped" -- Sections/04c-reencoding.tex, "Second, the scope is the
%     plate".
%   [DOWN] registered atomic / unregistered splits into sub-words --
%     Sections/03-Apparatus.tex line 283, and 00-HowToRead.tex line 209.
%
% WHAT IS ILLUSTRATIVE, AND SAYS SO IN THE CAPTION:
%   * the six gene profiles in stage 1. No axis is drawn; they carry no
%     measurement. Their ARITHMETIC is exact bar by bar, which is the only
%     thing a reader may take from them:
%       treated  {17, 9,14,12, 8, 6,11, 5}
%       control  {10,10, 8, 8, 7, 7,10, 6}
%       shift    { 7,-1, 6, 4, 1,-1, 1,-1}  = treated - control
%       generic  { 3, 2, 3, 2,-2, 2,-2, 2}
%       residual { 4,-3, 3, 2, 3,-3, 3,-3}  = shift - generic
%     Chosen so the residual carries the MAJORITY of the shift's squared norm
%     (74 of 106 here), because fig:residual-energy measures 62.1%
%     drug-specific and a figure drawing the residual as a sliver would
%     contradict it. Both figures are UNCENTRED sums of squares, each a
%     squared L2 norm and NOT a variance: the shift {7,-1,6,4,1,-1,1,-1} has
%     mean 2, not 0, and 49+1+36+16+1+1+1+1 = 106; the residual gives 74 the
%     same way. The label fig:residual-energy and the file
%     figs/04c-residual-energy.tex keep the older spelling because they are a
%     cross-reference target and an \input target.
%     v1 drew {1,0,0,1,0,1,0,0}: five of its eight bars had height zero and
%     rendered as nothing at all, so the one object the thesis is about read
%     as "missing" or "to be determined". FIXED HERE, and drawn SIGNED,
%     because the sign is what splits the two blocks in stage 2.
%   * which gene sits at which coordinate in stage 3, and which two the
%     prediction puts on the wrong side of the boundary. The STEM HEIGHTS are
%     not illustrative -- they are the exact weights.
%   * no cosine VALUE is printed anywhere. The two stem vectors are drawn, not
%     measured, so a number there would be an invented result.
%
% LAYOUT NOTES, all forced by a build:
%   * three stages, two dividers, one header rule, one brace. The brace spans
%     all THREE because the retrieval reference is an end-of-pipeline quantity:
%     it is a retrieval score computed in the scored space of stage 3, so
%     stopping it at stage 2 would misattribute it.
%   * stage 2's target string is ONE line. Wrapped over two it read as two
%     separate sequences and lost the up -> [DOWN] -> down ordering, which is
%     the only reason [DOWN] is drawn at all. Every chip is a NAMED node
%     chained off its predecessor's east anchor, so the two braces attach to
%     real node boundaries rather than to hand-computed x's. That was not
%     fastidiousness: the hand estimate of the string's width was 6mm short of
%     the truth, and on the first build [END_CELL] crossed the divider into
%     stage 3.
%   * type sizes follow figs/circularity.tex: \footnotesize (10pt, the caption
%     size) for the sentences a reader must read, \scriptsize (8pt) for symbol
%     labels, chips, brace labels and subordinate notes. Nothing is \tiny.
%   * stage 1 labels are centred on profiles 1.85cm apart. At 1.45cm
%     "generic(c)" and "residual r(d,c)" touch (v1's note, still true). They
%     are \scriptsize and not \footnotesize for the same reason: at 10pt they
%     collide even at 1.85cm, which a build showed.
% ===========================================================================
\begin{figure}[htbp]
\begin{fullwidth}
\centering
\begin{tikzpicture}[
  x=1cm, y=1cm, font=\small,
  bar/.style={fill=ink!72, draw=none},
  faintbar/.style={fill=ink!30, draw=none},
  base/.style={rulegrey, line width=0.25pt},
  gene/.style={draw=rulegrey, rounded corners=1pt, inner sep=1.6pt,
               font=\ttfamily\scriptsize, text=ink, anchor=west},
  shared/.style={gene, fill=rulegrey!12, text=quietink, draw=rulegrey!55},
  differs/.style={gene, fill=rulegrey!32, draw=ink!75, thick},
  atomic/.style={rounded corners=1pt, inner sep=1.6pt, fill=ink, draw=ink,
                 font=\ttfamily\scriptsize, text=white, anchor=west},
  dots/.style={anchor=west, font=\scriptsize, text=quietink},
  lbl/.style={font=\scriptsize, text=quietink},
  strong/.style={font=\scriptsize, text=ink},
  said/.style={anchor=west, font=\footnotesize, text=ink},
  op/.style={font=\normalsize, text=quietink},
  stagehead/.style={anchor=west, font=\sffamily\footnotesize\bfseries,
                    text=accent},
  note/.style={anchor=north west, font=\scriptsize, text=quietink,
               align=flush left, inner sep=0pt},
]

% ---- one bar profile. #1 x, #2 baseline y, #3 signed heights, #4 style.
% Heights are in "units"; one unit is 0.052cm. Bars are 0.125cm wide on a
% 0.160cm pitch, so eight span 1.245cm. The baseline rule is drawn under EVERY
% profile, because three of the six are signed and a bar below the line has to
% read as negative rather than as a misprint.
\newcommand{\fframeprof}[4]{%
  \foreach \hh [count=\ii from 0] in {#3} {%
    \pgfmathsetmacro{\bxx}{#1 + \ii*0.160}%
    \pgfmathsetmacro{\byy}{\hh*0.052}%
    \fill[#4] (\bxx,#2) rectangle ++(0.125,\byy);%
  }%
  \draw[base] (#1,#2) -- ++(1.245,0);%
}

% ---- one signed-rank stem plot. #1 zero-line y, #2 slot/weight list.
% Slots sit on a 0.255cm pitch from x = 13.16; a weight of 1 is 0.34cm.
\newcommand{\fframestems}[2]{%
  \draw[base] (13.10,#1) -- (17.02,#1);%
  \foreach \ss/\ww in {#2} {%
    \pgfmathsetmacro{\sxx}{13.16 + (\ss-1)*0.255}%
    \pgfmathsetmacro{\syy}{#1 + \ww*0.34}%
    \draw[ink!72, line width=0.55pt] (\sxx,#1) -- (\sxx,\syy);%
    \fill[ink!72] (\sxx,\syy) circle (0.5pt);%
  }%
  \draw[rulegrey, line width=0.3pt] (12.96,{#1-0.34}) -- (12.96,{#1+0.34});%
  \draw[rulegrey, line width=0.3pt] (12.92,{#1+0.34}) -- (13.00,{#1+0.34});%
  \draw[rulegrey, line width=0.3pt] (12.92,{#1-0.34}) -- (13.00,{#1-0.34});%
  \node[anchor=east, font=\scriptsize, text=quietink]
       at (12.88,{#1+0.34}) {$+1$};%
  \node[anchor=east, font=\scriptsize, text=quietink]
       at (12.88,{#1-0.34}) {$-1$};%
}

% ================================================== the header band
\node[stagehead] at (0.05,3.80) {1\,\textperiodcentered\,the residual is built};
\node[stagehead] at (5.70,3.80) {2\,\textperiodcentered\,it becomes a string};
\node[stagehead] at (12.60,3.80) {3\,\textperiodcentered\,it becomes a vector};
\draw[rulegrey, line width=0.4pt] (0.05,3.62) -- (17.30,3.62);
\draw[rulegrey!70, line width=0.3pt] (5.50,-1.66) -- (5.50,3.54);
\draw[rulegrey!70, line width=0.3pt] (12.40,-1.66) -- (12.40,3.54);

% ================================================== STAGE 1: two subtractions
\node[lbl]    at (0.67,3.36) {treated};
\node[lbl]    at (2.52,3.36) {control};
\node[strong] at (4.37,3.36) {shift $\shift(d,c)$};
\fframeprof{0.05}{2.30}{17,9,14,12,8,6,11,5}{bar}
\node[op] at (1.53,2.58) {$-$};
\fframeprof{1.90}{2.30}{10,10,8,8,7,7,10,6}{bar}
\node[op] at (3.38,2.58) {$=$};
\fframeprof{3.75}{2.30}{7,-1,6,4,1,-1,1,-1}{bar}

\node[strong] at (0.67,1.90) {shift};
\node[lbl]    at (2.52,1.90) {generic$(c)$};
\node[strong] at (4.37,1.90) {residual $\resid(d,c)$};
\fframeprof{0.05}{1.30}{7,-1,6,4,1,-1,1,-1}{bar}
\node[op] at (1.53,1.42) {$-$};
\fframeprof{1.90}{1.30}{3,2,3,2,-2,2,-2,2}{faintbar}
\node[op] at (3.38,1.42) {$=$};
\fframeprof{3.75}{1.30}{4,-3,3,2,3,-3,3,-3}{bar}

\node[note, text width=5.30cm] (nG) at (0.05,0.88)
  {generic$(c)$ is the mean over the \emph{other} training drugs in the same
   plate group; a group with fewer than three of them yields no generic and the
   condition is dropped. The residual is drawn signed because its sign is what
   splits the two blocks in stage~2.};

% ================================================== STAGE 2: the encoding fork
% Every gene label is a real panel gene, so every string drawn could occur in
% this data. The shared positions carry four of the most abundant panel genes by
% mean measured expression over the 606 consensus truth profiles
% (nir_consensus_profiles.npz: GAPDH rank 1, RPS6 2, HSPD1 7, ALDOA 8; GAPDH and
% RPS6 agree with gene_identity_check.json). They were ACTB, TPT1 and EEF1A1,
% none of which is in the 946-gene panel. In row (b) DDB2, an in-panel p53
% target, replaces MDM2, which is not in the panel. CDKN1A and HMOX1 are in the
% panel and unchanged. Labels only: no position, count or number changed.
\node[said] at (5.62,3.30) {\textbf{(a)} genes ranked by \emph{expression}};
\node[shared]  (a1) at (5.62,2.88) {GAPDH};
\node[shared]  (a2) at ([xshift=3pt]a1.east) {RPS6};
\node[shared]  (a3) at ([xshift=3pt]a2.east) {HSPD1};
\node[differs] (a4) at ([xshift=3pt]a3.east) {CDKN1A};
\node[shared]  (a5) at ([xshift=3pt]a4.east) {ALDOA};
\node[dots]    (a6) at ([xshift=4pt]a5.east) {$\cdots$};
\node[shared]  (b1) at (5.62,2.48) {GAPDH};
\node[shared]  (b2) at ([xshift=3pt]b1.east) {RPS6};
\node[shared]  (b3) at ([xshift=3pt]b2.east) {HSPD1};
\node[differs] (b4) at ([xshift=3pt]b3.east) {HMOX1};
\node[shared]  (b5) at ([xshift=3pt]b4.east) {ALDOA};
\node[dots]    (b6) at ([xshift=4pt]b5.east) {$\cdots$};
\node[said] at (5.62,2.06)
  {only \textbf{34.6 / 200} tokens differ \;(\textbf{17\%})};

\node[said] at (5.62,1.56)
  {\textbf{(b)} genes ranked by \emph{signed residual}};
\node[differs] (c1) at (5.62,1.14) {CDKN1A};
\node[differs] (c2) at ([xshift=3pt]c1.east) {DDB2};
\node[dots]    (c3) at ([xshift=3pt]c2.east) {$\cdots$};
\node[atomic]  (c4) at ([xshift=4pt]c3.east) {[DOWN]};
\node[differs] (c5) at ([xshift=4pt]c4.east) {HMOX1};
\node[dots]    (c6) at ([xshift=3pt]c5.east) {$\cdots$};
\node[atomic]  (c7) at ([xshift=4pt]c6.east) {[END\_CELL]};
\draw[decorate, decoration={brace, mirror, amplitude=3pt}, rulegrey]
  (c1.south west) ++(0,-0.06) -- ($(c3.south east)+(0,-0.06)$);
\node[lbl, anchor=north west] at ($(c1.south west)+(0,-0.16)$)
  {$100$ genes, $\resid > 0$};
\draw[decorate, decoration={brace, mirror, amplitude=3pt}, rulegrey]
  (c5.south west) ++(0,-0.06) -- ($(c6.south east)+(0,-0.06)$);
\node[lbl, anchor=north west] at ($(c5.south west)+(0,-0.16)$)
  {$100$ genes, $\resid < 0$};
\node[said] at (5.62,0.20)
  {\textbf{118.2 / 200} tokens differ \;(\textbf{59\%})};

\node[atomic] (dn) at (5.62,-0.28) {[DOWN]};
\node[note, text width=5.35cm] (nD) at ($(dn.east)+(0.14,0.11)$)
  {is \emph{one} registered atomic token. Unregistered it splits into
   sub-words, and the up\slash down boundary stops being a clean symbol.};

% ================================================== STAGE 3: the scored vector
\node[note, text width=4.70cm] (nW) at (12.60,3.44)
  {signed rank weights $w_r = 1/\log_2(r{+}1)$};

% the two coordinates on which the prediction and the truth carry OPPOSITE
% SIGNS -- the boundary error the [DOWN] token exists to prevent. Drawn FIRST
% so the stems sit on top of it.
\fill[rulegrey!18] (14.82,0.70) rectangle (15.33,2.42);
% The band's label sits in the GAP BETWEEN the two plots, running right from the
% band's right edge. Above the plots it ran into the second line of the weights
% note, which is math and taller than a text line; centred on the band it set
% flush against "true v" and the two read as one phrase. In the gap the nearest
% ink is the predicted plot's slot-10 stem, which stops at y = 1.72.
% anchor=EAST at 17.30, not west at 15.42: anchored west, this one 2.07cm label
% was the widest thing in the picture and pushed the bounding box out to
% 497.60pt, 1.35pt past the fullwidth measure -- one Overfull \hbox, which
% tools/build.sh caps at zero.
\node[lbl, anchor=east] at (17.30,1.46) {opposite signs};

\node[strong, anchor=west] at (13.16,2.50) {predicted $\hat{v}$};
\fframestems{2.06}{1/0.6309, 2/1.0000, 3/0.5000, 4/0.3869, 5/0.4307,
                   6/0.3562, 7/0.3155, 8/-0.6309, 9/0.3333, 10/-1.0000,
                   11/-0.5000, 12/-0.4307, 13/-0.3562, 14/-0.3869,
                   15/-0.3333, 16/-0.3155}

\node[strong, anchor=west] at (13.16,1.50) {true $v$};
\fframestems{1.06}{1/1.0000, 2/0.6309, 3/0.5000, 4/0.4307, 5/0.3869,
                   6/0.3562, 7/0.3333, 8/0.3155, 9/-1.0000, 10/-0.6309,
                   11/-0.5000, 12/-0.4307, 13/-0.3869, 14/-0.3562,
                   15/-0.3333, 16/-0.3155}

\node[note, text width=4.70cm] (nA) at (12.60,0.50)
  {$946$ coordinates, $200$ of them non-zero. $\cos(\hat{v},v)$, ranked
   against the other drugs in the same cell line, is $\NIR$.};

\node[note, text width=4.70cm] (nB) at (12.60,-0.60)
  {\DOWN{} has a token position in stage~2 but \emph{no coordinate} here: the
   sign carries the boundary.};

% ================================================== the invariant
\draw[decorate, decoration={brace, amplitude=4pt}, quietink, line width=0.5pt]
  (0.05,-1.80) -- (17.30,-1.80);
\node[anchor=north, font=\footnotesize, text=ink] (nBR) at (8.67,-1.98)
  {and the within-well split-half retrieval reference is \textbf{unchanged}:
   $0.742 \rightarrow 0.743$};

% ---- MEASUREMENT, left in and inert. \fframedebug is undefined in main.tex, so
% this expands to nothing there; figs/_harness_frame.tex \def's it and the
% picture then reports its own extent, the right edge of the target string, and
% the east edge of every node that could plausibly be the widest thing in the
% drawing. All three mattered: the string crossed a divider on the first build,
% and the bounding box ran 1.35pt past the fullwidth measure on a later one.
\ifx\fframedebug\undefined\else
  \path let \p1=(current bounding box.south west),
            \p2=(current bounding box.north east),
            \p3=(c7.east), \p4=(a6.east) in
    \pgfextra{\typeout{FRAME bbox x \x1 .. \x2 ; y \y1 .. \y2}%
              \typeout{FRAME target string ends \x3 ; row (a) ends \x4}};
  \path let \p5=(nW.east), \p6=(nA.east), \p7=(nB.east), \p8=(nBR.east),
            \p9=(nD.east), \p{10}=(nG.east) in
    \pgfextra{\typeout{FRAME east nW \x5 nA \x6 nB \x7 nBR \x8 nD \x9 nG \x{10}}};
\fi
\end{tikzpicture}
\caption[The residual, the target string and the scored vector]{\captiontitle{The
information was always present; the standard target's tokens barely expose it
--- and the string the model writes is not the vector that scores it.}
\emph{Stage 1.} The drug-specific residual is built by two subtractions. The
control removes the cell's own state, leaving the \emph{shift}; the mean over
the other drugs in the same plate group removes the \emph{generic} programme
that every drug triggers, leaving what makes this drug different. The six
gene profiles are illustrative shapes and carry no measurement --- no axis is
drawn --- but the arithmetic within the figure is exact bar by bar, and the
residual is drawn signed because its sign is what splits the two blocks in
stage~2. \emph{Stage 2.} That same measurement, written two ways. Ranking genes
by expression (a) leads every drug's profile with largely the same abundant
genes: only $34.6$ of the top $200$ tokens differ between two drugs in the same
context; ranking by signed residual
(b) raises that to $118.2$. (Divergence counts quoted at cell-line scope, as in
Table~\ref{tab:divergence}; plate scope is similar, at $120.3$.) Row (b) is the
target string itself: a $100$-gene up block, \DOWN, a $100$-gene down block,
\ENDCELL. \emph{Stage 3.} What $\NIR$ ranks is not that string but a signed
vector over the $946$ panel genes, top-$k$ up positive and top-$k$ down
negative. The stem heights are the exact weights $1/\log_2(\textrm{rank}+1)$, so
the eight drawn per sign are $1.000$, $0.631$, $0.500$, $0.431$, $0.387$,
$0.356$, $0.333$ and $0.316$; which gene sits at which coordinate is
illustrative, and no cosine value is printed, because these two vectors are
drawn rather than measured. Stage~2 and stage~3 are different objects and only
stage~3 is scored. The bracket is the point of the figure --- the within-well
split-half retrieval reference is essentially identical under both encodings, so
no information has been added. What changed is how much of it the token sequence
exposes to the loss --- a property of the encoding, which is not the same as
showing the encoding is what binds (Table~\ref{tab:divergence}).}
\label{fig:frame}
\end{fullwidth}
\end{figure}

  \paragraph{Reading a generation back, and what it is scored against}
  \Cref{sec:apparatus-vocab} gives the target's two-block form and the map $\phi$
  that turns a signed signature into a vector; \cref{fig:frame} draws the whole path,
  from the two subtractions to the vector that is scored. Three things are specific to
  scoring a \emph{generation} and are fixed here. A generation supplies its ranks by
  token order within each block, and a token outside the panel or already placed is
  skipped without consuming a rank, so a block fills to at most $100$ genes. The
  sampled generations of one condition are averaged as vectors and the mean is scored
  once, rather than scored one by one and averaged. And the score is $\NIR$ over the
  other drugs in the same cell line --- a set that spans plates, in apparent violation
  of the same-plate rule of \cref{sec:res-blind}, and one that stands only because the
  plate-scoped generic has already subtracted each plate's mean shift, in the
  expectation sense the scope limit above fences off.
  
  
  % PLACEMENT ONLY. claimbox is `breakable`, and it was breaking here: three lines
  % and the accent rule at the foot of one page, the remaining ten at the head of
  % the next, so the box read as two fragments and the verdict lost its frame.
  % \budgettick reserves 5\baselineskip, which is not enough for a claim this long;
  % 16 is the whole box (13 lines plus its padding and two rules).
  \needspace{16\baselineskip}%
\beatsection{R}{5}{A per-drug constant does not exhaust the target}{sec:res-transfer}
\label{sec:methods-vardecomp}

\begin{question}[5]
The residual says what the model should predict, not how much there is to
predict, and the same training result reads differently against a large reachable
signal and a small one. Two measurements settle the scale: what share of a
perturbation response is drug-specific at all, and how much of that share a table
of per-drug constants would already capture. The first has a clean answer. The
second has an honest one.
\end{question}

% NAMING. The quantity below is the squared Euclidean norm of the shift, and the
% prose says so. shared/inference.py :: energy_shares computes it with nothing
% centred -- denom = s @ s, then (r @ r)/denom, (g @ g)/denom and 2*(r @ g)/denom
% -- so it is a sum of squares about the origin and NOT a variance, whatever that
% function's "Exact variance decomposition" docstring says. Two names in this
% source keep the older signal-processing word on purpose and must not be
% renamed: \label{fig:residual-energy} is a cross-reference target (it prints as
% "Figure 4.12", so the word never reaches the page) and figs/04c-residual-energy
% is an \input target. The reader sees "squared norm" everywhere; only these two
% identifiers still read "energy".
\noindent
This section measures scale rather than skill. The residual says what the model
should predict; it says nothing about how much there is to predict, and the same
training result reads very differently against a large reachable signal and a
small one. Two measurements settle it, and this section takes them in order. The
first asks how much of a perturbation response is drug-specific at all. The second
asks how much of \emph{that} a table of per-drug constants would already reproduce
with no model involved.

The first is a decomposition, and it is available because the residual was defined
as a subtraction.  Since $\resid = \shift - \textrm{generic}$, the shift is the sum
of its two parts, and squaring gives
\begin{equation}
 \lVert \shift \rVert^{2}
 \;=\; \lVert \resid \rVert^{2}
 \;+\; \lVert \textrm{generic} \rVert^{2}
 \;+\; 2\,\langle \resid,\, \textrm{generic} \rangle .
\end{equation}
Dividing through by $\lVert \shift \rVert^{2}$ gives three shares that sum to one
by construction. Measured across the cache,
$\lVert \resid \rVert^{2} / \lVert \shift \rVert^{2} = 0.621$ (drug-specific),
$\lVert \textrm{generic} \rVert^{2} / \lVert \shift \rVert^{2} = 0.390$ (shared component), and the
cross term $2\langle \resid, \textrm{generic}\rangle / \lVert \shift \rVert^{2}
= -0.011$. The names on those three terms come from how the generic is built, not from the
algebra. \Cref{sec:apparatus-vocab} forms it as the mean shift of the \emph{other}
training drugs in the condition's own (cell line, plate) group, with the
condition's own drug struck out --- so it is an estimate of what a cell in that
context does when a drug is applied to it, constructed without reference to
\emph{which} drug. That is what makes $\lVert \textrm{generic} \rVert^{2}$ the part
of the response attributable to being treated at all, and leave-one-drug-out is
what guarantees it: were the drug's own conditions inside its own generic, part of
what that drug does would be counted as shared. The residual is then defined as
everything the shift carries that this estimate does not, so
$\lVert \resid \rVert^{2}$ is the part attributable to the identity of the drug. Both follow from the construction, and the measurement
decides only how the magnitude divides between them. 
The cross term is the substantive number, because it is the only one of the three
that could have come out otherwise, and what it reports is a geometry rather than
a share. Orthogonality between residual and generic is exactly the statement
$\langle \shift, \textrm{generic}\rangle = \lVert \textrm{generic} \rVert^{2}$,
which is the condition for the coefficient in a regression of the shift on the
generic to equal one. Measured here that coefficient is $0.986$. The generic is
therefore subtracted with the weight it actually carries: the plate-local,
leave-one-drug-out mean neither over- nor under-removes the shared programme, and
the residual retains no systematic trace of it. That is an empirical check on the
estimator of \cref{sec:methods-residual}, not a consequence of how it was defined.

Two readings follow. The first is that the drug-specific and shared components are
separate \emph{directions} rather than two settings of one. Had drugs differed
mainly in how strongly they provoke a common stress programme, the residual would
lie along the generic and the cross term would be large; at $-0.011$ it is not.
Drugs differ in the direction of their response, which is the property any
drug-conditional prediction must exploit and the one a scalar potency cannot
express. The second reading is size. With the components perpendicular the norms
compose by Pythagoras, $\lVert \resid \rVert$ and $\lVert \textrm{generic} \rVert$
standing at $0.786$ and $0.620$ of $\lVert \shift \rVert$, so the drug-specific
direction carries more magnitude than the shared one, roughly three to two, and a
shift sits about $38^{\circ}$ off its own drug-specific axis.

This is also what reconciles the decomposition with the counts of
\cref{tab:divergence}, where subtracting the control added $76.8$ differing
positions and subtracting the generic only $8.9$ more. The generic is
\emph{common mode}. Two drugs in one group are handed nearly the same generic,
differing only by which single drug leave-one-out withholds from a mean over the
group, so subtracting it translates both shifts by almost the same vector and
leaves the difference between them almost unchanged. A component can accordingly
carry two fifths of the squared norm and still barely move a count of between-drug
differing positions; the $8.9$ it does move is reranking, since a ranking is not
invariant under translation. Large magnitude and small contrast contribution are
not in tension, and the two measurements say different things about the same
object.

One qualification travels with all of it. What has been decomposed is the
\emph{shift}, treated minus control, whereas a cell sentence ranks the treated
profile itself, so no term of this decomposition appears anywhere in the standard
target (\cref{fig:residual-energy}).

% ===========================================================================
% 04c-residual-energy.tex -- the encoding problem in one picture.
% \input by Sections/04c-reencoding.tex. Compiles against thesis_v2/preamble.tex
% and nothing else; no external data, no \includegraphics.
%
% EVERY NUMBER IS SOURCED. Sibling 04c-residual-energy.json carries the key
% paths; the short form:
%   sq.norm  RESULTS_cluster/vardecomp_matched.json -> main.scope
%            residual_energy_share 0.6211205563374411 -> 62.1%
%            generic_energy_share  0.38993489697995193 -> 39.0%
%            cross_energy_share   -0.011055451267510967 -> -0.011
%            energy_shares_sum     1.000000002049882   -> 1.000
%            (the key names say "energy"; the quantity is the squared L2 norm)
%            (text, all three agreeing; line numbers drift, so each carries the
%             phrase to search for:
%             04c-reencoding.tex:301-303  "decomposes exactly into residual"
%             01-Introduction.tex:313-315 "Decomposing the treated-minus-control"
%             05-Limitations.tex:240-242  \paragraph{Frame coverage}
%             -- verify by the phrase, not the number, before quoting a line)
%   tokens   RESULTS_cluster/target_divergence.json, K=200
%            scopes.cellline.full.divergence_mean     34.594876910336524 -> 34.6 (17%)
%            scopes.cellline.shift.divergence_mean   111.4231960024091   -> 111.4 (56%)
%            scopes.cellline.residual.divergence_mean 118.24622167432057 -> 118.2 (59%)
%            (text: tab:divergence, 04c-reencoding.tex:42-44, the three body
%             rows of that table -- agrees)
%   the caption's 120.3 is scopes.plate.residual.divergence_mean 120.29533990815328
%   and is quoted in prose only; it is NOT drawn, because a plate-scope bar in a
%   cell-line-scope block is exactly the cross-block subtraction this layout is
%   built to prevent.
%
% LAYOUT NOTES, all forced by a build:
%   * two SUB-BLOCKS, one axis, one divider. The blocks are DIFFERENT QUANTITIES
%     (a squared-norm share and a count of token positions); the divider, the
%     1.68cm gap, the 2:1 bar heights (4.8mm vs 2.4mm) and the unit sentence
%     are what stop a reader subtracting across them. Nothing spans the divider.
%   * TOTAL WIDTH 14.191cm = 403.775pt against \textwidth = 404.02913pt (142.0mm,
%     measured in _v.log; preamble.tex:59 records the 132->142mm revision). The
%     figure fills 99.94% of the measure and needs no fullwidth, so the page it
%     lands on may still carry a margin note. It was previously laid out for the
%     132mm block, drew 348.3pt = 122.4mm and under-filled by 19.6mm.
%     Geometry: \X 4.900cm label gutter + \L 9.190cm per unit share; the right
%     end of the overhang is \X+1.011*\L = 14.19109cm. Picture box measured by
%     _scratch/figs/_harness_resenergy_w.tex, which reports
%     PICW = 403.77531pt PICH = 124.33833pt in _scratch/_harness_resenergy_w.log.
%   * ROW LABELS ARE ONE LINE EACH. Measured at \scriptsize in Pagella:
%     "residual --- the drug-specific target" 122.69pt = 4.312cm is now the
%     widest ("the perturbation response" measures 91.40pt), so the label column
%     stays at 4.65cm. Two-line labels at 0.62cm row pitch collide, which is why
%     the width is measured and not guessed.
%   * \bx must expand WITHOUT braces: it is used bare, as (\bx{0.5},y), and
%     inside an expression, as ({\bx{0.621}+0.14},y). A braced body makes the
%     second form "{{...}+0.14}" and pgfmath aborts with "Missing number".
%   * the 62.1% and 39.0% labels sit INSIDE/AFTER their segments; the token-block
%     labels sit AFTER the accent segment, because the 17% segment is 1.56cm
%     wide and will not hold its own label.
%   * THE HALF-WAY RULE IS DRAWN ONLY WITHIN THE TWO BAR BANDS. The previous
%     version ran it floor-to-ceiling and struck through both grey annotations
%     ("shift's" in the cross-term line, "positions" in the unit line); the old
%     header's guarantee that the rule "passes through none of them" guarded the
%     bold in-bar labels only and never covered the prose. Measured clearances
%     at bx{0.5} = 9.495cm, all in the bar bands: "62.1% drug-specific" ends at
%     7.50cm (1.99cm clear), "34.6 (17%)" ends at 8.04cm (1.46cm clear), and
%     "111.4 (56%)", "118.2 (59%)" and "generic 39.0%" all begin right of the
%     rule at 10.19, 10.46 and 10.75cm. The rule is drawn last, over the bars.
%     The two segments are given the same vertical extent as the upright ink
%     rule at 1.000, so the figure's two verticals rhyme instead of competing.
%   * THE "half" LABEL SITS UNDER THE TOKEN BLOCK, not at the top of the figure.
%     The half-way mark does its work below the divider, where 17% / 56% / 59%
%     are read against it; at the top it was as far from that reading as the
%     picture allowed.
%   * the generic segment runs on from where the residual segment ends, out to
%     1.011, with the ink rule at the shift's own squared norm: the
%     decomposition is exact, and drawing the -0.011 overhang rather than
%     describing it keeps the two shares from being silently renormalised to
%     sum to one.
%   * the cross-term annotation is anchored EAST and grows leftwards. Anchored
%     WEST at the right end of the bar it walked 51.1pt off the page.
% ===========================================================================
\begin{figure}[htbp]
% MARGINALIA (06-figures.md): wider than the 124mm text column, so it
% sits in fullwidth; the caption below spans the text column only.
\begin{fullwidth}\raggedright
\begin{tikzpicture}[x=1cm, y=1cm, font=\small]
  \def\X{4.900}                  % left edge of every bar
  \def\L{9.190}                  % bar length for 100%
  \newcommand{\bx}[1]{\X+#1*\L}  % value -> x   (no braces: see header)
  % PIN THE BOUNDING BOX. Left free, the picture measures 409.59pt and runs
  % 5.56pt over the block: the rowname nodes are anchor=east at 4.75 with a
  % 4.65cm text width, so their inner sep hangs ~0.09cm left of x=0, and the
  % "half" descender hangs below. Neither carries ink -- the row labels are
  % flush right and set single-line, so their ink starts at x=0.07cm. Pinning
  % to [0, \bx{1.011}] therefore clips nothing and makes the drawn width equal
  % the stated 14.191cm = 403.775pt exactly.
  \useasboundingbox (0,-0.25) rectangle ({\bx{1.011}},4.12);
  \tikzset{
    rowname/.style={anchor=east, text width=4.65cm, align=flush right,
                    font=\scriptsize, text=quietink},
    inlab/.style={anchor=west, font=\scriptsize\bfseries, text=white},
    outlab/.style={anchor=west, font=\scriptsize\bfseries, text=accent},
    greylab/.style={anchor=west, font=\scriptsize, text=ink}}

  % ========== SUB-BLOCK 1: the response, by squared norm. Bars 4.8mm. ======residuak
  \node[rowname] at (4.75,3.70) {the perturbation response};
  \fill[accent]      (\X,3.46)          rectangle (\bx{0.621},3.94);
  \fill[rulegrey!45] (\bx{0.621},3.46)  rectangle (\bx{1.011},3.94);
  \node[inlab]   at ({\X+0.14},3.70)          {$62.1\%$ drug-specific};
  \node[greylab] at ({\bx{0.621}+0.14},3.70)  {generic $39.0\%$};

  % the shift's own squared norm, and the cross term as the overhang past it
  \draw[ink, line width=0.7pt] (\bx{1.0},3.32) -- (\bx{1.0},4.08);
  \node[anchor=north east, text width=7.6cm, align=flush right,
        font=\scriptsize, text=quietink] at ({\bx{1.011}},3.28)
    {the upright rule marks the shift's own squared norm; the overhang past it
     is the cross term $-0.011$};

  % --- the divider: the two sub-blocks are different quantities ------------
  \draw[rulegrey, line width=0.4pt] (0,2.44) -- ({\bx{1.011}},2.44);
  \node[anchor=north west, font=\scriptsize\itshape, text=quietink] at (0,2.38)
    {above: shares of the shift's squared norm \quad below: shares of $200$
     token positions};

  % ========== SUB-BLOCK 2: the target, by token position. Bars 2.4mm. ======
  % r1 -- the standard target
  \fill[rulegrey!25] (\X,1.54) rectangle (\bx{1.0},1.78);
  \fill[accent]      (\X,1.54) rectangle (\bx{0.17},1.78);
  \node[rowname] at (4.75,1.66) {full profile --- the standard target};
  \node[outlab]  at ({\bx{0.17}+0.14},1.66) {$34.6$ \;($17\%$)};

  % r2 -- the shift
  \fill[rulegrey!25] (\X,0.92) rectangle (\bx{1.0},1.16);
  \fill[accent]      (\X,0.92) rectangle (\bx{0.56},1.16);
  \node[rowname] at (4.75,1.04) {shift (treated $-$ control)};
  \node[outlab]  at ({\bx{0.56}+0.14},1.04) {$111.4$ \;($56\%$)};

  % r3 -- the residual
  \fill[rulegrey!25] (\X,0.30) rectangle (\bx{1.0},0.54);
  \fill[accent]      (\X,0.30) rectangle (\bx{0.59},0.54);
  \node[rowname] at (4.75,0.42) {residual - (shift $-$ generic)};
  \node[outlab]  at ({\bx{0.59}+0.14},0.42) {$118.2$ \;($59\%$)};

  % --- the half-way rule, drawn LAST so the bars do not bury it, and drawn
  % ONLY within the two bar bands so it crosses no prose. See header.
  \draw[ink, line width=0.5pt, densely dotted] (\bx{0.5},3.32) -- (\bx{0.5},4.08);
  \draw[ink, line width=0.5pt, densely dotted] (\bx{0.5},0.18) -- (\bx{0.5},1.90);
  \node[anchor=north, font=\scriptsize\itshape, text=ink] at (\bx{0.5},0.14) {half};
\end{tikzpicture}
\caption[The drug is most of the response, little of the
target]{\captiontitle{The drug is the majority of the perturbation response and almost
none of the standard target's tokens.} Above the rule: the shift's squared
norm decomposes exactly into a drug-specific residual, a generic programme and
a cross term, so $62.1\%$ of what a perturbation does to a cell is specific to
the drug. Below it: of the $200$ positions a cell sentence spends on the
standard target, only $34.6$ carry a different gene for two drugs in the same
context --- the other $83\%$ are the same string. Ranking the same measurement
by signed residual raises that to $118.2$, and at the plate scope this build
trains in the residual row reads $120.3$ ($60\%$). The two halves of the figure
are different quantities, a squared-norm share and a count of token positions;
they share the axis and the half-way rule and nothing else, and no difference
between them is a number. What the picture states is an ordering: the signal is
a majority of the response, the standard encoding exposes a small minority of
it to the loss, and the residual encoding puts the exposed fraction back on the
majority side (\cref{tab:divergence}).}
\label{fig:residual-energy}
\end{fullwidth}
\end{figure}



\paragraph{The second question is whether that majority has to be predicted at all}
The first measurement says the drug-specific part is the larger half of a
perturbation response. It says nothing about whether that part is hard to produce.
A residual could be large and still be a per-drug \emph{constant} --- the same
signature for a given drug wherever it is applied --- in which case it is
reproducible from a table indexed by drug name, with no model, no prompt and no
cell involved. So the second question is how much of the residual is that
constant, and how much of it moves with the context the drug is applied in.

Writing
\begin{equation}
 \resid(d,c) \;=\; \beta(d) \;+\; \kappa(d,c) \;+\; \epsilon,
 \qquad
 \Tcoef \;=\; \frac{\operatorname{var}\beta}{\operatorname{var}\beta + \operatorname{var}\kappa},
\end{equation}
splits the residual into a per-drug constant, a context-specific remainder and
noise. $\beta(d)$ is the drug \emph{main effect} and is what a lookup table
reproduces with no model at all;\seealso{\cref{sec:res-lookup} builds that lookup
and scores it.} $\kappa(d,c)$ is the only component a conditional model could win
on; and $\Tcoef$ is the share of the structured part that the constant carries.
The baseline this sets up is not invented for the occasion. It is DrEval's naive
predictor~\cite{dreval} in transcriptomic space: their
\texttt{NaiveMeanEffectsPredictor} predicts $\hat{y}_{ij} = \mu^c_i + \mu^d_j -
\mu$, cell-line and drug main effects with no interaction term, and
$\textrm{generic}(c) + \beta(d)$ is that quantity lifted from a scalar potency to a
$946$-dimensional profile. Asking what $\kappa$ holds is therefore asking what a
model could add over the field's own no-skill predictor.

\paragraph{The direct estimator is unavailable, and what replaces it}
Fitting the variance components would be the obvious route: estimate
$\operatorname{var}\beta$ and $\operatorname{var}\kappa$ and read $\Tcoef$ off
them. It cannot be done on this cache. Separating $\operatorname{var}\kappa$ from
$\operatorname{var}\epsilon$ needs two independent measurements of the same drug in
the same context, and $\num{286}$ of the $\num{287}$ distinct treatments here occur
in exactly one well, so the model is unidentified on its own data. The route taken
instead buys identification from a quantity already computed for the reliability
line. Two residuals of one drug in two cell lines are compared by cosine; that
cosine is depressed by measurement noise in both, and dividing by the geometric
mean of the two conditions' split-half reliabilities removes the depression.

% PLACEMENT: this box carries a display equation and cannot break usefully. At
% the previous setting it opened at the foot of the page and tcolorbox broke it
% immediately after the title bar, stranding "DEFINITION The transfer
% coefficient" alone at the page bottom with every line of the body overleaf.
% \budgettick's own \needspace{5\baselineskip} is not enough to catch that; the
% box is ~13 lines tall, so ask for that and let it move whole.
\needspace{12\baselineskip}
\begin{defn}{The transfer coefficient}
With $\num{286}$ of $\num{287}$ treatments in exactly one well there is
no replicate from which to separate $\operatorname{var}\kappa$ from
$\operatorname{var}\epsilon$, so the variance-components model is unidentified on
its own. The cosine route buys identification through split-half reliability:
\begin{equation}
 \hat{\Tcoef} \;=\; \frac{\cos\!\big(\resid(d,c_1),\, \resid(d,c_2)\big)}
 {\sqrt{\rho(d,c_1)\,\rho(d,c_2)}},
 \qquad
 \rho(d,c) \;=\; \frac{2\cos(\resid_A, \resid_B)}{1 + \cos(\resid_A, \resid_B)},
\end{equation}
where $\resid_A$ and $\resid_B$ are disjoint halves of condition $(d,c)$ and
$\rho(d,c)$ is the Spearman--Brown reliability of that condition's residual at
its full cell count.
\end{defn}

We name explicitly two terms in the formula. The
\emph{cross-line cosine} is the numerator alone: $\cos(\resid(d,c_1), \resid(d,c_2))$,
the plain cosine similarity between one drug's two full residuals in two different
cell lines, computed with nothing else applied to it. \emph{Disattenuation} is a
correction, not a description, and it targets a specific bias. Each residual
entering that cosine is itself a noisy estimate --- it carries the same sampling
error the reliability line was built to detect --- and noise in two vectors being
compared pulls their cosine toward zero independently of whatever the vectors
would agree on if measured exactly. So the raw cross-line cosine understates the
true agreement between contexts by an amount set by how noisy each residual is,
and it understates it more where reliability is lower. Dividing by
$\sqrt{\rho(d,c_1)\,\rho(d,c_2)}$, the geometric mean of the two conditions' own
split-half reliabilities, removes exactly that shrinkage: this is the classical
correction for attenuation, applied here to a cosine instead of the correlation
coefficient it was derived for.\seealso{\cref{sec:appendix}'s planted-truth check
licenses using it on a cosine} $\hat{\Tcoef}$ is what the formula returns after
that division, and calling it \emph{disattenuated} names the correction it has had
applied, not a new quantity being introduced.

Whether the correction actually removes the shrinkage, rather than just
relabelling it, is checkable, because a case exists where the true answer is
known. Data were simulated with $\kappa = 0$ built in --- every context gets the
same residual for a given drug, so by the definition of $\Tcoef$ the truth is
exactly $1$. Measurement noise is simulated too, so the raw cross-line cosine on
this synthetic data still comes out below $1$: it reads $0.734$, which a reader
taking it at face value would misread as $27\%$ context-specific structure,
$1 - 0.734$, when by construction there is none. Dividing that same $0.734$ by
the geometric-mean reliability recovers $\hat{\Tcoef} = 1.001$
(\cref{tab:tcontrols}), against a truth of $1$. The correction is therefore shown
to work on a case where the answer is known, not merely assumed to work because
the derivation is standard, and every number below rests on that check.

\paragraph{What the coefficient is, and what it is not} Four qualifications attach
before any estimate is quoted, and together they make $\Tcoef$ a descriptive
sensitivity analysis rather than a fundamental measurement. First, $\kappa$ is
written as a drug\,$\times$\,cell-line interaction because the decomposition
assumes one, but the estimator does not identify such a term; $\kappa$ is therefore
called the \emph{context-specific remainder} and $\Tcoef$ a disattenuated
cross-context residual-similarity coefficient, never a variance share. Second, it might be tempting to read $1 - \Tcoef$ as the headroom above the constant baseline, however we advise against that: the cross-line pairs
entering it may differ in dose and in treatment well, so it mixes several
axes into one number. Third, the scoped calculation retains conditions at a
reproducibility threshold of $0.2$ and takes each residual against a
\emph{self-including} generic its own drug helped define, which is not the training
frame. Fourth, the reliability in the denominator is a within-well split-half
precision, optimistic relative to biological reproducibility --- and that optimism
has a direction. An over-large $\rho$ makes $\hat{\Tcoef}$ too small, so
$1 - \hat{\Tcoef}$ errs \emph{high} and bounds from above what a per-drug constant
leaves on the table, without measuring it.

One further caution is methodological rather than interpretive. Spearman--Brown is
derived for correlations between parallel halves, not for cosines, and what
licenses its use here is the planted-truth validation of \cref{sec:appendix}.


% ...........................................................................
% ...........................................................................
\begin{table}[htbp]
  % MARGINALIA: in-column; caption and the note in the margin.
  \begin{lrbox}{\tvfb}
  \begin{tabularx}{\textwidth}{@{}Y r l@{}}
  \thead{Quantity} & \thead{Value} & \\
  \midrule
  ground truth, by construction ($\kappa = 0$) & $\Tcoef = 1.000$ & \\
  \addlinespace[10pt]
  uncorrected cross-line cosine, $\cos(\resid(d,c_1), \resid(d,c_2))$ & $0.734$ & \\
  that cosine read naively as an interaction share, $1 - 0.734$ & $27\%$ & (truth:
    $0\%$) \\
  \addlinespace[10pt]
  disattenuated estimator, $\hat{\Tcoef}$ & \key{$1.001$} & \\
  \end{tabularx}
  \end{lrbox}
  \sidecap{table}{\captiontitle{The disattenuation correction, checked where the truth is
  known.} Residuals were simulated with $\kappa = 0$ built in, so by the
  definition of $\Tcoef$ the ground truth is exactly $1$; every row is computed on
  that synthetic data, not on the atlas.\label{tab:tcontrols}
  \par\addvspace{6pt}Measurement noise alone, with no cross-context difference anywhere in the
  simulation, is enough to pull two identical residuals' cosine down to $0.734$.
  The correction recovers the true value to within $0.001$ and licenses the
  reference column of \cref{tab:transfer}.}
  \end{table}

% ...........................................................................
\begin{table}[htbp]
% MARGINALIA: in-column; caption in the margin.
\begin{lrbox}{\tvfb}
\begin{tabularx}{\textwidth}{@{}Z{1.12} r l Z{0.88}@{}}
\thead{Quantity} & \thead{Estimate} & \thead{$95\%$ interval} & \thead{Reference} \\
\midrule
$\Tcoef$ (cross-context, plate-scoped generic) & $0.553$ & $[0.514, 0.593]$
  & simulated $\kappa{=}0$ null $= 1.001$ \\
similarity scale NOT shared across contexts, $1-\Tcoef$ & $44.7\%$
  & $[40.7\%, 48.6\%]$ & mixes cell line, dose, well and estimator \\
\addlinespace[10pt]
$\Tcoef$ across dose (same drug, same line) & $0.707$ & $[0.632, 0.782]$
  & reference scale \\
\addlinespace[10pt]
structure-matched negative control & $-0.008$ & spanning zero
  & must be $\approx 0$ \\
\end{tabularx}
\end{lrbox}
\sidecap{table}{\captiontitle{The transfer coefficient, estimated on real molar doses with
dyadic cluster-robust intervals.} The matched negative control sits on zero. The
same-plate control that produced the earlier false alarm reads $+0.485$ under a
cell-line-scoped generic. $\Tcoef$ is a cross-context similarity coefficient, not
a variance share, and $1 - \Tcoef$ must not be read as a
drug\,$\times$\,cell-line interaction share.\label{tab:transfer}}
\end{table}
\paragraph{What the estimate says} Two things about how this number was obtained
are worth restating before it is read. \emph{Real molar doses}: the dose entering
this calculation is the resolved concentration of \cref{sec:methods-data},
recovered from \texttt{drugname\_drugconc}.
\emph{Dyadic cluster-robust}: $\hat{\Tcoef}$ is computed over \emph{pairs} of
conditions, and two pairs sharing either a drug or a cell line are not
independent observations, so the interval is not a plain bootstrap over pairs
but the Fafchamps--Gubert dyadic estimator of \cref{sec:apparatus-clusters},
built for exactly that dependence.

Under that estimator, on real doses, $\hat{\Tcoef} = 0.553$ $[0.514, 0.593]$
(\cref{tab:transfer}). Read with every qualification above in force: a little over
half of the cross-context similarity scale is carried by the per-drug constant,
and $44.7\%$ $[40.7\%, 48.6\%]$ of it is not.

That reading rests on the negative control holding, so the control is worth
describing rather than only quoting. It is built by keeping everything the
estimand depends on fixed --- the same pair of cell lines, the same plate-scoped
generic, the same disattenuation --- and breaking only the one thing $\hat{\Tcoef}$
is supposed to measure: which drug is compared. Two \emph{different} drugs' residuals
are compared across the same cell-line pair instead of one drug's two residuals,
so any structure the control still detects would have to come from the pipeline
itself, the generic construction or the dyadic clustering, and not from drug
identity. It returns $-0.008$, an interval spanning zero. Nothing in the
construction manufactures apparent cross-context agreement on its own; what
$\hat{\Tcoef} = 0.553$ measures is therefore attributable to the shared drug, not
to an artifact of how the estimator is built.

There is consequently a conditional target that a per-drug constant does not
express, and its size is bounded from above at roughly $45\%$ of a similarity
scale --- not a variance share, and not a quantity any model has been shown to
reach.

\paragraph{Dose is the milder perturbation} The estimand $\hat{\Tcoef}$ is meant
to say how much of the context-specific remainder survives a change of
\emph{cell line} specifically, but the cross-context pairs computing it also
differ in well and, since doses are not matched across cell lines, in dose ---
so the $44.7\%$ that does not transfer could in principle be about that mixture
rather than about cell-line context as such. Dose supplies a cleaner comparison,
because it is the one axis that can be varied while everything else the
estimand depends on --- same drug, same cell line, same generic, same
disattenuation --- stays fixed. If a drug's residual moved as much between two
of its own doses as it does between two cell lines, cell line would be doing no
special work; if it moved much less, dose would be shown to be the milder
perturbation, and the part of the cross-context estimate that dose alone could
explain away would be small.

Holding drug and cell line fixed and varying only the molar concentration,
$\hat{\Tcoef} = 0.707$ $[0.632, 0.782]$ on $\num{363}$ pairs, against
$\hat{\Tcoef} = 0.553$ $[0.514, 0.593]$ pooled across cell lines; the intervals
do not overlap. Dose is therefore the milder perturbation: a drug's residual
survives a change of concentration better than it survives a change of cell
line, so the $44.7\%$ that does not transfer across context is not merely an
artefact of comparing different doses to each other. The comparison is not
perfectly clean --- the cross-context arm still differs in well and dose as well
as cell line, so this is the closest the section comes to a single-axis test
rather than one --- but the direction runs the way the estimand's own logic
predicts.


 
\paragraph{Whether anything could learn the remainder} The $44.7\%$ a per-drug
constant does not reach is not automatically a target only a full model could
hit. A simpler explanation is available first. If a cell line modulated every
drug's response in a broadly consistent direction --- a ``cell-line
personality'' that is not any drug's own signal but is nonetheless learnable
from that line's other conditions --- then two \emph{different} drugs measured
in the \emph{same} cell line should share more of that direction with each
other than either shares with a drug measured in a \emph{different} cell line,
once each drug's own main effect is removed. That prediction is specific enough
to test directly.

Testing it means removing the main effect first. Write
$\kappa(d,c) = \resid(d,c) - \hat{\beta}(d)$, where $\hat\beta(d)$ is estimated
\emph{leave-one-cell-line-out}: averaged over every cell line except $c$, so
condition $(d,c)$ never enters its own baseline. That exclusion is not a
formality. Estimated over every line including $c$, $\hat\beta(d)$ would
contain $\resid(d,c)$ itself, and centring a value on a mean that contains its
own contribution induces a $-1/(m-1)$ correlation between the centred values by
the arithmetic of the mean alone, before any real structure enters ---
precisely the leave-one-out concern already argued for the generic in
\cref{sec:methods-residual}, protecting a variance test here rather than a
training target. With the main effect removed, the test compares two cosines:
$\cos(\kappa(d,c), \kappa(d',c))$, two different drugs' remainders in the same
cell line, against $\cos(\kappa(d,c), \kappa(d',c'))$, two different drugs in
two different cell lines. A cell-line personality predicts the first running
higher than the second. Both cosines are disattenuated by $\kappa$'s own
split-half reliability before the comparison is made, for the same reason
$\hat{\Tcoef}$ itself is: an uncorrected cosine would read the personality's
absence as smaller than it is, purely because $\kappa$ is a noisier quantity
than $\resid$.

Averaged over lines, neither arm shows the signature the hypothesis predicts.
Under the plate-scoped generic the same-cell-line arm reads
\ci{-0.0149}{-0.0251}{-0.0050} and the different-cell-line arm
\ci{+0.0003}{-0.0071}{+0.0081}, an excess within line of $-0.0153$ for which no
interval was computed; both come from a cell-line-clustered bootstrap,
$\num{38}$ and $\num{42}$ clusters respectively, over $\num{4000}$ same-line and
$\num{3884}$ cross-line pairs drawn from $\num{1282}$ conditions. A real
per-line direction predicts the same-line arm sitting \emph{above} the
cross-line one; instead it sits below it, and it is the only one of the two
distinguishable from zero at all. That negative same-line reading is not itself
read as evidence of some opposite structure, because it has a known
non-biological cause: same-line pairs that also share a plate share a generic,
and subtracting one shared quantity from both members of a pair pushes their
cosine negative by construction --- the identical mechanism the different-drug
same-plate diagnostic reads at $-0.017$ elsewhere in this section. Cross-line
pairs share no generic and are unaffected by it. So the uncontaminated arm is
the cross-line one, and it is indistinguishable from zero, while the same-line
arm's departure from zero is explained by the construction rather than by any
drug-independent cell-line effect. Neither arm, then, shows the positive
within-line excess a consistent per-line direction would require. 


% ...........................................................................
\beatsection{R}{5}{Re-encoding produces prompt sensitivity}{sec:res-encoding}

\begin{question}[5]
Every target change tried so far has not shown robust counterfactual predictions. Does training on the residual change that? The instrument is the
stratified scramble, the same prompt with another drug's name and mechanism
substituted, read as a ladder of strata because the swapped-in drug is itself a
drug the model knows and so is not a neutral comparison. The model moves. It is
not the only target change that makes it move.
\end{question}

\noindent

Every scramble so far has been run on \texttt{tier2}: held-out drugs, ones the
checkpoint never saw in training, with the substituted partner drawn from that
same held-out set. That is a hard test and deliberately so, but it asks for two
things at once. A model that passes it must respond to the drug it is named
\emph{and} have learned something transferable about a drug it has never
encountered. Nothing has passed it. The single-cell arm is null under the
random-partner comparison and shows no ordering under its own designed sweep;
consensus is flat at every swap distance; optimal transport clears zero only
against the most dissimilar partners and nowhere before them.

This section separates the two demands and asks only the first. The ablation is
unchanged. Same prompt, same substitution, same comparison set, but it is
run on conditions the residual checkpoint actually trained on,and, for the two
controls, on a sample whose split its evaluation did not record --- so that for at
least the arm this section is about, nothing has to be extrapolated and
memorisation is available. That makes it the easier question by
construction, and the right one to ask next: if a model is sensitive to the drug
token at all, this is where it shows, and a null here would be much harder to
attribute to the difficulty of the split.

It also brings in a fourth target. The residual re-encoding of
\cref{sec:methods-residual} was not among the three constructions of
\cref{sec:res-epsilon}; it is what this section is about, and two of those three
are carried forward with it --- but not as a yardstick. Relaxing the split is a
concession, and the obvious objection to anything found after making it is that
the concession \emph{is} the finding: that scoring a model on conditions it has
already fitted is by itself enough to open a gap against a scrambled prompt. The
controls are here to answer that objection, not to compete for a ranking. They are
read out of the same evaluation cache, through a ladder built by the same rule,
against the same kind of comparison. One thing about them has to be said straight
away, because it bounds what they can settle: their evaluations carry no split
manifest, so which of their conditions each checkpoint had already seen is not
recorded, and they cannot certify that the relaxation was granted to them on the
same terms. What they do bound is the instrument --- whether this ladder, read out
of this cache, returns a large positive gap as a matter of course --- and that is
the reason they are drawn. No ordering between arms is claimed anywhere below.

Carrying them forward also fixes where the earlier nulls do \emph{not} come from,
and the two controls answer that differently. \Cref{sec:res-blind} and
\cref{sec:res-epsilon} found neither construction able to order its gap by swap
distance on held-out drugs, scored in the expression frame they were trained to
emit --- flat throughout for single cell, and for optimal transport a single
narrow excursion at the far end and nothing before it. Read here on seen conditions and
in the residual frame, the single-cell arm is still not recovered: its series spans zero
at every stratum. Neither control's split is recorded, so that is not the clean
seen-condition null it would be worth having; what it does establish is narrower
and still useful: pointed at this checkpoint, on this cache and this ladder, the
instrument returns nothing at all, so it does not open a gap merely by being run. Optimal
transport is partly recovered, and that is the second reason to draw it. Its
\texttt{orth} and \texttt{opposite} gaps both clear zero here, small but real, and
this section is where the reading that that arm does nothing is withdrawn. The
withdrawal is budget-driven and internal to this instrument: on the same
checkpoint against the same \texttt{orth} comparator, $600$ tokens give
\ci{+0.0144}{-0.0032}{+0.0323}, spanning zero, and $1400$ give
\ci{+0.0292}{+0.0144}{+0.0441}, which does not --- though on a pool sharing only
$\num{10}$ conditions with the first. It is not a re-run of
\cref{sec:res-epsilon}'s own sweep, which is a different frame, population and
comparator; what it withdraws is that sweep's conclusion, not its number. What it still does not show is ordering across the
ladder, and the limits below are explicit that this is a statement
about the ladder its pool supports rather than about the arm. Neither reading
turns on a comparison with the residual arm, and
\cref{sec:res-generalisation} is where the harder question goes back to whichever
arm survives this one: whether the response holds on conditions never trained
on.

One consequence of that relaxation belongs up front, before any design detail or
any number: the arms are no longer scored on a common set. Each is read on a
scoring sample of a few hundred conditions rather than a full population, since
three scrambled generations per condition is expensive, and those samples were
drawn separately --- three pools in all: one $\num{250}$-condition pool, shared by
all three arms at $600$ tokens and reused for the pre-repair checkpoint at
$1400$; the repaired residual arm's $\num{300}$ at $1400$; and the controls'
$\num{248}$ at $1400$. At $1400$ the two are mostly non-overlapping
and they differ in size. What is held fixed is not the
data but the shape of the instrument: the near/orth/opposite ladder below is
built the same way wherever it is built, so the only reading any of it supports
is an arm against its own scramble on its own conditions, never one arm's number
against another's.

The instrument itself needs one design choice made explicit before any number is
worth reading. For each condition, a partner drug is drawn from every other
condition sharing its cell line \emph{and} its plate but naming a different drug
--- same plate so a swap never differs from the target in batch as well as in
identity, different drug so a drug's own other dose is never mistaken for a
scramble --- and three fixed partners are taken from that pool by measured
residual cosine: the \emph{near} partner is the most similar available,
\emph{orth} the one closest to orthogonal, \emph{opposite} the most
anti-correlated. The pool is ranked by residual cosine and not by the
\texttt{moa-fine} mechanism annotation because the annotation does not track it:
two drugs sharing a mechanism label sit a mean distance $2.322$ apart against
$2.377$ for two that do not ($\num{1244}$ same-mechanism against $\num{17563}$
different-mechanism pairs, from the drug atlas of \cref{sec:apparatus-vocab}), just over two per cent closer,
so ``different mechanism'' does not imply ``different response''. A condition
with fewer than three eligible partners is dropped rather than assigned a
contaminated one. The three partners are fixed by this construction and reused, unchanged, across
every arm scored within one evaluation run --- though not across runs: the
$1400$-token control run draws on a different pool, whose \texttt{opposite} rung
is much the weaker of the two, as \cref{fig:repair} sets out.

On a scoring sample of $\num{300}$ conditions the holdout assigns to
\texttt{train}, $\num{293}$ of which the evaluation also flags as seen --- the
\texttt{train} tier of the three-way holdout \cref{sec:methods-residual} builds,
scored on that same checkpoint's evaluation (\texttt{re\_v3.json}, one
generation budget of $1400$ tokens throughout, the only budget this checkpoint
has been scored at) --- the model separates from its own scramble on two of the
three partners, and moves in the predicted direction on the third, in the order
the construction predicts throughout. Against \texttt{near}, the hardest test,
the gain is smallest and does not clear zero,
\ci{+0.0229}{-0.0232}{+0.0689}; against \texttt{orth} it is
\ci{+0.0898}{+0.0382}{+0.1415}, clustered on cell line and well
(\ci{+0.0898}{+0.0291}{+0.1506} clustered on drug and well); against
\texttt{opposite} it more than doubles again,
\ci{+0.1942}{+0.1358}{+0.2527}.\gloss{scramble}{The same prompt with the drug's
name and mechanism replaced by another drug's, in three strata by residual
similarity.} One checkpoint, one evaluation, one generation budget, one
population: on the conditions this checkpoint was trained on, the model separates
from its own scramble in the order the construction predicts, and it is not the
only target construction that separates at all.

 % ...........................................................................
% ...........................................................................
\paragraph{Why three strata, and why one of them is quoted} The three strata are
not three attempts at the same measurement. They are three detection
sensitivities, set deliberately. Every prediction is scored against the
\emph{real} drug's truth, so what changes across the strata is not how hard the
prediction is to produce but how visible it is that the model followed the drug
it was named. Substitute a near-twin and an unchanged prediction is nearly the
correct answer anyway, so the scramble scores well and there is almost nothing
to see. Substitute an opposite effect drug and a model that actually follows its
prompt is pushed away from the truth in the one direction the metric measures,
so the scramble collapses and the separation is at its widest. The model's own
score is fixed at $0.598$ throughout; it is the scramble NIR across the stratum that falls, $0.575$
to $0.508$ to $0.404$, and the gap that opens with it. A model doing genuine
counterfactual work should therefore separate most on \texttt{opposite} and
least on \texttt{near}, which is the ordering \cref{tab:reencode-gradient}
reports. That ordering is the evidence. A metric artefact, or a model
reproducing a generic response regardless of the drug, has no reason to arrange
itself by how dissimilar the substituted drug happens to be, so the shape of the
series across strata says something no single stratum can.

A single number is nonetheless needed, and not for anything internal to this
section: everything downstream requires one. Reporting an arm compactly means one figure per
arm, not three; \cref{sec:res-generalisation} compares the trained tier against
two held-out ones, and its evaluation was run against one fixed comparator
(\texttt{scramble\_orth}) rather than re-run three times; and every later section
that cites this effect --- the lookup-table comparison, the channel gate, the
chapter summary --- inherits whatever one quantity is named here. The choice is
therefore about which stratum's gap can honestly carry the label ``the effect''.

\texttt{Orth} is the only candidate, and the reason is that it alone has a
defensible null. Scored on its own, the \texttt{near} partner reaches $0.575$,
above the $0.500$ chance line: naming a similar drug is partly \emph{correct},
so a gap measured against it subtracts real signal and understates the effect.
The \texttt{opposite} partner scores $0.404$, below chance: that substitution is
actively wrong, so a gap measured against it inflates. \texttt{Orth} sits at
$0.508$, on the chance line almost exactly, which is what makes its gap read as
distance above chance rather than distance from a contaminated baseline. Reading
the three partners as points on one line rather than as three separate
comparisons makes the same point: their own scores fall monotonically with the
cosine between the substituted drug's truth and the real one, and \texttt{orth}
is where that line crosses chance. It is the one place on the ladder where
subtracting the comparator neither gives the model credit for a partly right
answer nor charges it for a wrong one.

% ...........................................................................

% ...........................................................................
% ...........................................................................
% ...........................................................................
\paragraph{How the scoring sample is built for each arm, and what makes the
three comparable} 
Nothing here is scored on a full population. For each arm, a
\emph{scoring sample} is drawn from the evaluation cache, for the residual arm
from conditions it had already seen during its own training --- a few hundred conditions, not the whole set --- and every
number in this section is read off that sample. The residual arm's is the
$\num{300}$-condition \texttt{train} tier of the three-way holdout
\cref{sec:methods-residual} builds, and its trained status is recorded per
condition. The two controls are read on one shared sample at each budget --- $\num{250}$
conditions at $600$ tokens and a near-disjoint $\num{248}$ at $1400$, with only
$\num{10}$ conditions common to the two --- drawn from the same evaluation cache
the residual arm is itself scored out of. Neither control evaluation carries a
split manifest; both record the split as unknown for every condition, so which of
them each checkpoint had already trained on is not established. Their rows are
nulls of unverified split rather than demonstrated seen-condition nulls, and the
seen-condition framing of this section is secure only for the residual arm, whose
$\num{300}$ come with the three-way holdout attached.

The three samples are not three views of one population, and this is the point
worth being plain about: the three arms were trained on different data. Each
target construction came with its own corpus, built at its own time, under its
own inclusion rules, and no arm was retrained to match another's. A scoring
sample drawn this way therefore means a different set of conditions
for each arm, and the samples differ in size as well: the residual arm's
$\num{300}$ and the controls' $\num{248}$ have just $\num{15}$ conditions in
common. Nothing survives that in the form a cross-arm comparison would need. What
each sample does support is one arm read against its own scramble on its own
conditions, and that is the only reading taken below. 

% ---------------------------------------------------------------------------
\begin{table}[htbp]
% MARGINALIA: full width, like the other interval tables of this section;
% caption above at the text width, notes below in the sans.
\caption{\captiontitle{On the conditions the checkpoint trained on, the gap runs from
a partly-correct partner to an actively wrong one, at one checkpoint and one
generation budget.} Model $\NIR$ is $0.598$ in every row; only the comparator
changes. \texttt{re\_v3.json} / \texttt{scramble\_stratum\_audit\_v4.json},
$n = \num{300}$ throughout.}
\label{tab:reencode-gradient}
\begin{fullwidth}
\begin{threeparttable}
\begin{tabular}{@{}l l r r r l@{}}
\thead{Stratum} & \thead{Partner is} & \thead{Partner cos.} & \thead{Scramble NIR}
  & \thead{Gap} & \thead{$95\%$ interval (cell line, well)} \\
\midrule
\arm{near} & most similar available & $+0.248$ & $0.575$ & $+0.023$
  & $[-0.023, +0.069]$ \\
\arm{orth} & closest to orthogonal & $-0.0001$ & $0.508$ & \key{$+0.090$}
  & $[+0.038, +0.142]$ \\
\arm{opposite} & most anti-correlated & $-0.261$ & $0.404$ & $+0.194$
  & $[+0.136, +0.253]$ \\
\end{tabular}
\begin{tablenotes}[flushleft]
\item Partner cosine is the mean cosine, in the residual frame, between a
condition's own residual and the partner substituted for that row's stratum. Scramble NIR is the NIR of the prediction when swapping the name of the drug with that of what is in the stratum(near,orth,opposite).
The \texttt{orth} row additionally excludes zero clustered on drug and well,
\ci{+0.0898}{+0.0291}{+0.1506}; no drug-clustered interval was computed for the
other two rows.
\end{tablenotes}
\end{threeparttable}
\end{fullwidth}
\end{table}

% .....................................................................
% RESOLVED IN PROSE (not in data): re_singlecell_1400.json has cache_dir=ot_cache,
% which is NOT the corpus pythia_sft_endcell trained on, so its rows are not a
% verified seen-condition test. Rather than quote numbers under a sampling
% discipline they do not have, the prose above now states this limitation
% outright and marks the single-cell rows as a null of unverified split. If the
% same-corpus rerun on data_diverse2_endcell_big ever lands, swap the numbers in
% that paragraph, the two-panel figure/JSON and ctrl1400_summary.json, and the
% caveat can be deleted.
% Verified 2026-08-30 from the artifacts: re_singlecell_1400.json and
% re_ot_1400.json score the SAME 248 conditions (split='unknown', trained=False,
% train_file=None, holdout=None); those 248 overlap the residual train tier in
% only 15 conditions; repro_cos min 0.2022 / median 0.2784 for the 248 against
% 0.1089 / 0.1827 for the residual 300 --- i.e. the controls' pool is the more
% reliable one, so the reliability filter is NOT a confound favouring the
% residual arm.

% ...........................................................................
\paragraph{Two controls, read at both budgets, and one residual cell that is a
different model}
 Every run this section has is in \cref{tab:budget-arms}, and the residual arm
contributes three of them rather than two. Two are the \emph{pre-repair}
checkpoint --- trained on targets built before the split-before-fit correction of
\cref{sec:methods-residual}, when the generic and the reliability filter were
still fitted over every condition and the holdout assigned afterwards --- scored
at both budgets on one $\num{250}$-condition pool. Those two rows are the
chapter's one clean budget comparison: one checkpoint, one pool, one comparator,
with the budget the only thing that moves, and the gap roughly doubling at every
stratum when it does. The third is the repaired checkpoint at $1400$ tokens, which
is where this chapter's canonical result comes from and which has never been
scored at $600$; it cannot be read against either pre-repair row without
confounding a target rebuild with a budget change. The two
control arms have no such problem at the level of the checkpoint: each was scored
at both budgets on its own single checkpoint. Their two budgets are not read on
one condition set, though --- the $600$- and $1400$-token pools share only
$\num{10}$ conditions --- so those rows move budget and population together, and
are not budget comparisons alone.

Read at $1400$ tokens, where the residual arm's canonical checkpoint actually
lives, neither control shows a three-stratum pattern --- though on their pool the
three strata are not three distinct tests, so only the near-versus-rest half of
that is a real null. The
single-cell control is flat and null at every stratum. The optimal-transport
control is not flat, and clears zero on two of the three strata --- but at
\texttt{orth} and \texttt{opposite} almost equally, which on this pool is close to
being the same measurement twice: in $\num{138}$ of its $\num{248}$ conditions the
two strata name the same partner drug. Both arms'
full three-stratum results at both budgets are in \cref{tab:budget-arms}.

One asymmetry in scrutiny attaches to the optimal-transport \texttt{orth} value.
Unlike the residual arm's own headline number, it has not been re-checked under
drug-and-well clustering; the drugs behind those conditions repeat across them
the same way the residual arm's do, so a wider drug-clustered interval is not
ruled out, only untried.

Three limits stop any of this being read as an ordering between arms. The first
is the frame. The residual checkpoint's generations are scored where they are
emitted, while the two controls emit cell sentences, whose decoded profile is a
function of gene rank order alone; a drug effect can reach their gaps only by
reordering top-ranked genes, which is the channel \cref{tab:divergence} measures
as nearly closed --- $34.6$ of $200$ top positions differing between two drugs
under the standard encoding, against $120.3$ under the residual one. Their gaps
are attenuated by construction, and an attenuated gap cannot be read against a
native one.

The second is that the ladder is not equally built in the two pools, and this is
measurable rather than a matter of interpretation. Across the residual arm's
$\num{300}$ conditions, every \texttt{opposite} partner is a different drug from
that condition's \texttt{orth} partner, and all $\num{300}$ opposite cosines are
negative, mean $-0.261$. Across the controls' $\num{248}$, \texttt{opposite} and
\texttt{orth} name the \emph{same} drug in $\num{138}$ of them --- $\num{56}\%$
--- and only $\num{133}$ of the opposite cosines are negative at all, mean
$-0.009$. The controls' top two rungs are largely one rung, which is the plainest
reading of why their \texttt{orth} and \texttt{opposite} gaps come out almost
equal, and it means an absence of near-to-opposite ordering in those arms is not
evidence that they lack it: the instrument that would have shown it was not there
to be failed.

The third is that these are separate evaluations over different, mostly
non-overlapping condition sets drawn from separately built corpora, not a paired
comparison, and the arms were not held to the same training discipline.

One thing that could have been a further limit is not one, and it is worth saying
so rather than leaving it to be assumed. The residual arm's $\num{300}$ conditions
are the reliability-surviving population of \cref{sec:methods-residual}, which
invites the objection that its result is a property of a pre-screened sample. The
controls' $\num{248}$ are the more reliable pool of the two, not the less: their
lowest reproducibility cosine is $+0.202$ against the residual tier's $+0.109$,
their median $+0.278$ against $+0.183$. Whatever this section's numbers are, they
are not the product of the residual arm having been handed the easier
conditions.

A pattern that survives the sharpest stratum on one arm is evidence about that
arm. Its absence on an arm whose sharpest stratum was never built is not evidence
about that one, and neither fact is a ranking. The comparative case in this
chapter is not made in this section at all: it is made by \cref{tab:divergence},
which counts how far two drugs' target strings diverge under each encoding and
needs no model to do it.
\begin{figure}[htbp]
\begin{fullwidth}
  \raggedright
  \includegraphics{fig-repair}
  \caption[Budget, not target, moved this gap]{\captiontitle{Raising the generation
  budget from $600$ to $1400$ tokens sharpened a real pattern rather than
  manufacturing one: the residual checkpoint's advantage over its own scramble
  grew at every stratum, the optimal-transport control's grew from
  null-adjacent to clear, and the single-cell control stayed flat at both.}
  Residual frame throughout; the optimal-transport and single-cell series are the
  \emph{implied} residual described above, since neither checkpoint was trained
  to emit one. Both panels share the residual checkpoint's superseded, pre-repair
  run, the only one with a $600$-token measurement; the current checkpoint's
  canonical result (\cref{tab:reencode-gradient}) is not drawn here. The $x$-axis
  is the stratum by name, not by measured cosine: at $600$ tokens all three series
  share one $250$-condition partner pool (mean residual cosine $+0.406$,
  $-0.0016$, $-0.329$ for near, orth, opposite); at $1400$ tokens the residual
  series keeps that same pool while the two controls were scored on a different,
  smaller $248$-condition pool whose near and orth partners shift by less,
  $+0.403$ and $+0.036$, but whose \texttt{opposite} partner is only weakly
  anti-correlated, $-0.009$, so the two panels' \texttt{opposite} points are not a
  matched pair. Intervals are $95\%$ bootstrap, $n = 250$ per point at $600$
  tokens and $n = 248$ at $1400$ for the two controls; clustering differs by
  series and is not uniform across this figure --- the two controls' $1400$-token
  intervals are clustered on cell line and well, matching the number quoted for
  them in the text, while every other series here (the residual arm at both
  budgets, both controls at $600$) has only a coarser cell-line-only interval
  computed and is plotted with that instead.}
  \label{fig:repair}
\end{fullwidth}
\end{figure}
% ---------------------------------------------------------------------------
\begin{table}[htbp]
  % MARGINALIA: full width; caption above at the text width, notes below.
  % Every interval sits in an l column so it never breaks mid-cell.
  \caption{\captiontitle{Every arm at both generation budgets, model minus its own
  stratum-matched scramble.} The four control rows hold one checkpoint each across
  the two budgets, but not one condition set: the $600$- and $1400$-token pools
  share only $\num{10}$ conditions. The two pre-repair residual rows are the
  cleanest pair here --- one checkpoint, one $\num{250}$-condition pool, budget the
  only thing that moves. The repaired row is a different checkpoint from those two,
  so reading it against them would mix a target rebuild with a budget change.}
  \label{tab:budget-arms}
  \begin{fullwidth}
  \begin{threeparttable}
  \begin{tabular}{@{}l r r l l l@{}}
  \thead{Arm} & \thead{Budget} & \thead{$n$}
    & \thead{near} & \thead{orth} & \thead{opposite} \\
  \midrule
  \arm{residual, repaired}\tnote{a} & $1400$ & \num{300}
    & \ci{+0.023}{-0.023}{+0.069}
    & \key{\ci{+0.090}{+0.038}{+0.142}}
    & \ci{+0.194}{+0.136}{+0.253} \\
  \arm{residual, pre-repair}\tnote{b} & $1400$ & \num{250}
    & \ci{+0.035}{+0.006}{+0.066}
    & \ci{+0.072}{+0.037}{+0.106}
    & \ci{+0.143}{+0.111}{+0.179} \\
  \arm{residual, pre-repair}\tnote{b} & $600$ & \num{250}
    & \ci{-0.005}{-0.040}{+0.028}
    & \ci{+0.035}{-0.005}{+0.077}
    & \ci{+0.072}{+0.032}{+0.110} \\
  \addlinespace[10pt]
  \arm{optimal transport} & $1400$ & \num{248}
    & \ci{+0.009}{-0.005}{+0.024}
    & \ci{+0.029}{+0.014}{+0.044}
    & \ci{+0.028}{+0.013}{+0.044} \\
  \arm{optimal transport} & $600$ & \num{250}
    & \ci{+0.007}{-0.008}{+0.023}
    & \ci{+0.014}{-0.003}{+0.032}
    & \ci{+0.026}{+0.007}{+0.048} \\
  \addlinespace[10pt]
  \arm{single cell} & $1400$ & \num{248}
    & \ci{-0.021}{-0.059}{+0.016}
    & \ci{-0.013}{-0.046}{+0.020}
    & \ci{-0.022}{-0.049}{+0.004} \\
  \arm{single cell} & $600$ & \num{250}
    & \ci{-0.014}{-0.045}{+0.018}
    & \ci{-0.021}{-0.050}{+0.006}
    & \ci{-0.020}{-0.054}{+0.010} \\
  \end{tabular}
  \begin{tablenotes}[flushleft]
  \item[a] The canonical result, repeated from \cref{tab:reencode-gradient} for
  comparison. It is the only row scored on a population with a split manifest
  attached, and the only one the chapter's claim rests on.
  \item[b] Pre-repair checkpoint, trained on targets built before the
  split-before-fit correction of \cref{sec:methods-residual}, scored at both
  budgets on one $\num{250}$-condition pool. These two rows are the series
  \cref{fig:repair} draws for the residual arm. The repaired checkpoint has never
  been scored at $600$ tokens, and is not drawn in that figure.
  \end{tablenotes}
  \end{threeparttable}
  \end{fullwidth}
  \end{table}

  \paragraph{The slope carries its own null}
  One further check supports reading \texttt{orth} as the honest point on that
ladder, and it does so without selecting any comparator at all. Each condition
contributes three scrambled prompts, and each names a partner drug whose true
residual sits at a measured cosine from the real drug's --- on this tier
$+0.25$, $0$ and $-0.26$ for near, orth and opposite.  Regressing each prompt's $\NIR$ on
that cosine, over all three prompts of every condition at once, asks a question
none of the individual gaps can: does the score track how similar the named drug
actually is? The regression is pooled least squares with one intercept, three
points per condition and no condition-level intercepts absorbed, and the
target-drug prompt is never a point on the line --- only the three scrambles
are, so nothing in the fit contrasts the model against a chosen baseline.

Zero is the null by construction. A model that ignores the drug it is told
emits the same prediction whichever partner is substituted, so its three points
sit at the same height and the line is flat, whatever the partners' cosines
happen to be. On the trained tier the slope is \ci{+0.3728}{+0.2084}{+0.5139}:
positive, excluding zero, and steep enough that moving a partner from
anti-correlated to similar --- a span of roughly $0.5$ in cosine --- buys about
$0.19$ in $\NIR$, which is the same order as the spread between the three
comparator scores the ladder is built from. What that establishes is not
\texttt{orth} directly but the ladder underneath it: the comparators really do
fall in the order their labels claim, and by an amount that scales with the
cosine, so \texttt{orth}'s position at the chance crossing is a property of the
measured geometry and not an artifact of which three partners happened to be
picked. \Cref{sec:res-generalisation} fits the same slope on the two held-out
tiers. One caveat on reading its magnitude rather than its sign: with a single
pooled intercept the zero null is exact within a condition, but the fitted slope
also carries a between-condition component.


\paragraph{Both budgets, side by side} The budget is not a neutral runtime
setting here. Raising it from $600$ to $1400$ tokens moved the
optimal-transport control's \texttt{orth} gap from spanning zero to excluding
it, very nearly doubling the point estimate on an identical checkpoint and an
identically constructed comparator, though on a nearly disjoint sample of
conditions, while leaving the single-cell control null at both budgets --- so
the budget sharpens a gap that is already there and does not manufacture one
where there is none. \Cref{fig:repair} puts all three arms at both budgets side
by side, one panel each, across the three strata;
\cref{tab:budget-arms} carries the same numbers.

The figure's $x$-axis is the stratum by name rather than measured cosine,
because the two panels do not rest on one partner pool. At $600$ tokens all
three arms share a single $\num{250}$-condition pool. At $1400$ the two controls
were re-run on a smaller $\num{248}$-condition pool whose \texttt{opposite}
partner is barely anti-correlated, $-0.009$ against $-0.329$ in the pool the
residual series keeps; near and orth move far less, $+0.406\rightarrow+0.403$
and $-0.0016\rightarrow+0.036$. The controls' $1400$-token \texttt{opposite}
point is therefore a weaker test than any other point on either panel, and
plotting all three series on one continuous cosine axis would draw a shared
difficulty ladder the data does not support.


% ...........................................................................
\beatsection{R}{6}{Responding to the drug is not predicting it}{sec:res-generalisation}
\label{sec:methods-holdout}

\begin{question}[6]
A model that moves when the drug token moves has cleared a low bar, because
reproducing a training average requires that much too. Everything scored so far
came from conditions the model trained on, so what has been shown is that the
token is consulted, not that anything transferable was learned. Separating the
two needs a holdout, and the
atlas's own randomisation decides which one means anything: whole treatment
wells. Prompt sensitivity survives that test. Whether the residual itself
transfers comes back unsettled.
\end{question}

\noindent

\Cref{sec:res-encoding} established one thing, and deliberately only one. On
conditions the checkpoint had trained on, re-encoding the target made the model
respond to the drug it was named: the separation opens in the order the
construction predicts, it holds clustered on the cell line and on the drug, and
the comparator-free slope excludes zero. That was the easy question by design.
Nothing had to be extrapolated, memorisation was available, and a model that did
no more than reproduce a training average would have cleared much of the same bar.
What it establishes is that the prompt is consulted, not that anything was learned
that travels.

Where the response stops is the question worth asking, and this atlas supports two
ways of asking it. One recombines what the model has already met: a trained drug
seen again in a treatment well held out whole. The other withholds the drug from
fine-tuning altogether. The first is the weaker test and the second is the one the
claim has to survive. Both are reported below, and they fail in different places.

The physical layout of \cref{sec:methods-data} forces the choice. A (drug, cell
line) holdout, the field's standard Leave-Pairs-Out split and the one an earlier
version of this work used, leaves the held-out condition's treated well sitting
in the training set through its other $\approx 48$ cell lines. Measured against
that holdout, $\num{124}$ of the $\num{255}$ treated wells for which split
information exists placed a held-out condition and a training condition in the
same well.\seealso{\cref{sec:res-lookup}: what a held-out gap has to
beat.} That estimand measures within-atlas matrix completion, not independent
generalisation, and no processing of the targets changes it.

Holding out whole wells is the only assignment-respecting choice the design
offers, and it answers a narrower question, generalisation to an unseen treatment
well of a seen drug. Fixed-dose cross-context transfer, present in one line and
absent in another, is structurally unidentifiable in this atlas.

\paragraph{What each split holds} \texttt{train} takes every condition of the
remaining wells, $\num{4999}$ before the reliability line and 2,685
after it.
\texttt{unseen\_combo} holds $\num{36}$ treated wells out whole ($\num{904}$
conditions, $\num{900}$ after inventory filters), and it never takes a drug's
last training well; without that guard a two-well drug could silently become an
unseen drug while sitting in the unseen-combination arm, and the two arms would
measure the same thing. \texttt{unseen\_drug} holds every condition of
$\num{10}$ held-out drugs ($\num{714}$ conditions, $\num{711}$ after inventory
filters). Zero wells contribute conditions to both sides.

\texttt{unseen\_drug} is a control and not a third result. Tahoe fine-tuning
never sees the token, so a large effect there would indict the comparator, not
report a biological finding. ``Unseen'' means held out from fine-tuning with
mechanism supplied, not unknown to the pipeline; the Pythia-1b backbone was
pretrained on natural language and the prompt still carries the
\texttt{Mechanism:} field.

Two warnings ride with the design. DrEval reports, citing~\cite{partin}, that
apparently good performance on pair-holdout splits can be reached simply by
predicting each drug's average response across the training cell lines, which is
why the baseline ladder of \cref{sec:res-lookup} is not optional. And the
evaluation draws under a per-split quota, because a uniform draw reproduces the
pool's proportions and starves the split the experiment exists to test; in a
first run,
$\num{250}$ constructed held-out combinations yielded only $\num{33}$ scored
conditions. Under the quota the evaluation scores $\num{1394}$ conditions in all,
$\num{300}$ trained and $\num{199}$ unseen-drug, the rest unseen combinations,
almost the whole pool of $\num{900}$ that arm affords, while
\texttt{unseen\_drug} remains a quota of the $\num{711}$ available.

% ...........................................................................
\paragraph{The first result of this design is about the comparator} Which partner
defines the gap decides the size of every number in this section, and only one of
the three has a claim to the role. A gap splits at chance into how far the model
rises above it when handed the right drug and how far the comparator falls below
it when handed a wrong one,
\begin{equation}
 \gap \;=\; \underbrace{\big(\NIR_{\textrm{model}} - 0.5\big)}_{\textrm{how much naming the truth helps}}
 \;+\; \underbrace{\big(0.5 - \NIR_{\textrm{scramble}}\big)}_{\textrm{how much naming the lie hurts}},
\end{equation}
and only \texttt{orth} makes the second term vanish: its own score is $0.501$,
against $0.550$ for \texttt{near} and $0.423$ for \texttt{opposite}.

That second term is not noise --- it is the model being pushed away from a truth
it was told to miss --- but how far it is pushed depends on which partner the
prompt named. Each condition's three partners are drawn from the other drugs on
its own plate and ranked by the cosine between their true residual and the
target's (\cref{sec:res-encoding}): \texttt{orth} is the partner closest to zero,
and \texttt{opposite} the most anti-correlated available, meaning the drug whose
measured response points most nearly against the target's. That cosine is a fact
about which drugs shared a plate, not about the model, and it does not hold still
across the splits. Averaged over conditions, the \texttt{opposite} partner reaches
$-0.2615$ on \texttt{train} but only $-0.1842$ on \texttt{unseen\_combo}, because
the most anti-correlated drug \emph{available} is whatever happened to be plated
alongside; \texttt{orth} sits within $0.0005$ of zero on all three splits, since a
partner near zero is one every pool can supply.

That is what decides the choice. This section's claim is a comparison
\emph{between} splits --- whether a separation found on trained conditions
survives on held-out ones --- so the comparator has to mean the same thing in
every row, or part of any change from one row to the next is the plate rather than
the model. On the numbers above, \texttt{opposite} fails that condition and
\texttt{orth} meets it. Every gap in \cref{tab:generalisation} is therefore taken
against \texttt{orth}, and what is reported is the difference itself,
$\NIR_{\textrm{model}} - \NIR_{\textrm{orth}}$, never one term of the split
above.
% ...........................................................................
\begin{table}[htbp]
% MARGINALIA: full width; the two gap columns are interval strips on one axis
% (-0.12 to +0.16, hairline at zero). Caption above at the text width, notes
% below in the sans. Every number is the one the previous build printed.
\caption{\captiontitle{Whole-well holdout on the corrected evaluation, with the truth
verified against the build's own fit digest and the full held-out population
scored.} The two right-hand columns carry the SAME point estimate and differ only
in what they treat as the unit of generalisation. On \texttt{unseen\_combo} the
gap excludes zero clustered on cell line and does not clustered on drug, which is
why that row is reported as \notestablished. The \texttt{unseen\_drug} row
corroborates but does not prove: its \texttt{orth}-partner gap is negative while
its \texttt{opposite} gap is $+0.0379$. The strips share one axis, $-0.12$ to
$+0.16$, with a hairline at zero; a filled point excludes zero, a hollow one does not.}
\label{tab:generalisation}
\begin{fullwidth}
\begin{threeparttable}
\begin{tabular}{@{}l l l l l@{}}
\thead{Split} & \thead{$n$ (lines; drugs; wells)} & \thead{model $\NIR$}
  & \thead{$\gap$, clustered on cell line}
  & \thead{the same, clustered on drug} \\
\midrule
\arm{train} & $300$ ($39$; $76$; $136$) & $0.598$ \quiet{$[0.549, 0.647]$}
  & \ivstrip{0.0382}{0.1415}{0.0898}{$+0.0898\;[+0.0382,\,+0.1415]$}
  & \ivstrip{0.0291}{0.1506}{0.0898}{$+0.0898\;[+0.0291,\,+0.1506]$} \\
\addlinespace[10pt]
{\color{accent}\arm{unseen_combo}} & \needsaudit{895} ($42$; $34$; $36$)
  & $0.533$ \quiet{$[0.492, 0.575]$}
  & \ivstrip{0.0069}{0.0594}{0.0332}{$+0.0332\;[+0.0069,\,+0.0594]$}
  & \ivstrip[neg]{-0.0031}{0.0695}{0.0332}{$+0.0332\;[-0.0031,\,+0.0695]$} \\
\addlinespace[10pt]
\arm{unseen_drug} & $199$ ($40$; $10$; $28$) & $0.459$ \quiet{$[0.404, 0.514]$}
  & \ivstrip[neg]{-0.0836}{0.0206}{-0.0315}{$-0.0315\;[-0.0836,\,+0.0206]$}
  & \ivstrip[neg]{-0.1049}{0.0419}{-0.0315}{$-0.0315\;[-0.1049,\,+0.0419]$} \\
\end{tabular}
\begin{tablenotes}[flushleft]
\item Every gap printed here is measured against the geometry-defined
\texttt{orth} partner, whose own score is empirically near chance on these
conditions; \texttt{opposite} figures never appear as a column. Model $\NIR$
intervals test against the chance value $0.5$ with no comparator. Intervals are
two-way cluster-robust over cell lines and wells, with a Student $t$ reference
on $\min(G_{\textrm{line}}, G_{\textrm{well}}) - 1$ degrees of freedom, which is
why the well counts are printed. The \texttt{unseen\_drug} arm spans $28$
wells, and $t_{27} = 2.05$ against $1.96$ widens its interval by five per cent
(\cref{sec:methods-data}).
\end{tablenotes}
\end{threeparttable}
\end{fullwidth}
\end{table}

\paragraph{What the two held-out rows report} The gap is a paired contrast --- the
same checkpoint on the same conditions, differing only in the drug the prompt
names --- so it can be informative where an absolute score is not. On
\texttt{unseen\_combo} naming the right drug does help, by $+0.0332$: clustered on
cell line and well the interval excludes zero, \ci{+0.0332}{+0.0069}{+0.0594}, but
clustered on drug and well it does not, \ci{+0.0332}{-0.0031}{+0.0695}, and drug
is the axis a claim about transferring a \emph{drug}-specific residual has to
survive. That first interval is also the tighter of the two measurements
\cref{tab:generalisation} affords on this row --- a half-width of $\pm 0.026$
against the $\pm 0.042$ on the model's own score against chance quoted below,
both two-way cluster-robust on cell line and well --- because differencing against
the scramble removes what the two prompts share and buys a precision the absolute
score cannot have. Read that way, naming the right drug does help there, by
$+0.0332$: clustered on cell line and well the interval excludes zero,
\ci{+0.0332}{+0.0069}{+0.0594}, but clustered on drug and well it does not,
\ci{+0.0332}{-0.0031}{+0.0695}, and drug is the axis a claim about transferring a
\emph{drug}-specific residual has to survive. Positive for $20$ of $34$ drugs at a
sign test of $p = 0.39$, the effect is a majority tilt and not a uniform one, so
transfer to held-out combinations is not convincingly establsihed. On \texttt{unseen\_drug}
there is no benefit to the name at all: the gap is
\ci{-0.0315}{-0.0836}{+0.0206}, a point estimate favouring the \emph{scrambled}
prompt, and though the interval covers zero the direction is the one a model with
no knowledge of the held-out drug would give.

A second limitation sits beside that one and is not the same. The model's own
$\NIR$ is $0.533$ $[0.492, 0.575]$ on \texttt{unseen\_combo} and $0.459$
$[0.404, 0.514]$ on \texttt{unseen\_drug}, both covering chance: on a well it did
not train on, the prediction is no closer to that condition's truth than to the
truths of the drugs it is ranked against. Responding to the name and predicting
the response are separate properties, and this table finds the second absent on
held-out data whatever it finds of the first. Whether the first survives is asked
below, on an instrument that needs no comparator and reaches all three splits.

\begin{figure}[htbp]
\begin{fullwidth}
\raggedright
\includegraphics{fig-splits}
\caption[Chance-clearing on trained conditions only]{\captiontitle{The residual-frame
model clears chance on trained conditions only; the held-out-combination gap
excludes zero clustered on cell line but not clustered on drug, and the
unseen-drug gap flips sign.} Residual frame throughout, and $\NIR$ ranks each
prediction against the other drugs measured in the same cell line.
\textbf{(a)}~Left, model $\NIR$ against the chance value $0.5$ (solid rule), with
$95\%$ two-way cluster-robust intervals; $n = 300$, \needsaudit{895} and $199$ conditions on
\texttt{train}, \texttt{unseen\_combo} and \texttt{unseen\_drug}. The grey
dash on each row is that split's raw within-well split-half precision
reference, a measurement-precision figure computed from two halves of the same
well and not a replicate experiment; it carries no interval and is drawn as a
bare marker, and only \texttt{train} is filtered at the reliability line. Right,
on its own axis, the gap against the geometry-defined \texttt{orth} partner. Each
row has one point estimate, drawn twice with two intervals because the verdict
depends on which unit of generalisation is clustered: cell-line-clustered above,
drug-clustered below, each labelled with the Student~$t$ degrees of freedom it is
carried on. The grey band on the \texttt{unseen\_drug} row is that arm's
$80\%$-power detection limit around chance, $\pm 0.079$, the smallest departure
the arm could detect and not a confidence region. \textbf{(b)}~The per-condition
$\NIR$ distributions behind those three means, as reflected-boundary kernel
densities on a shared axis, each with its mean marked and the region at or below
chance shaded. The pooled $0.537$ is a mean over nearly flat distributions rather
than a location: even on \texttt{train}, $38.7\%$ of conditions sit at or below
chance, against $44.6\%$ and $55.8\%$ on the two held-out splits.}
\label{fig:splits}
\end{fullwidth}
\end{figure}


% FLAGGED, NOT FIXED -- the two panel-B numbers below (train $0.955$ and the
% spread $0.006$) are inherited from a superseded reconstruction taken at a
% +0.109 reliability line. This document defines the line at +0.1086: see :214
% ("$+0.109$, resolving to $+0.1086$"), :222-223 ("none falls at or below the
% resolved $+0.1086$") and :986 ("499 conditions sit below the $+0.1086$ line").
% Recomputed from RESULTS_cluster/re_v3.json -> records[] at repro_cos > 0.1086:
%   train         kept 300  0.9555006 -> 0.956   (:222-223 requires all 300)
%   unseen_combo  kept 396  0.9605642 -> 0.961   (:986's 499-below requires 396)
%   unseen_drug   kept 111  0.9597734 -> 0.960
%   spread: raw 0.0050636 -> 0.005, and max(rounded) - min(rounded)
%           = 0.961 - 0.956 = 0.005. BOTH conventions give 0.005 at the resolved
%           line, so this is NOT the convention clash it was reported as. 0.006
%           is just the old line's rounded-endpoint spread, 0.961 - 0.955. At
%           0.109 the kept counts are 299 / 394 / 111, which contradict :222-223
%           and :986 directly.
% The $0.132$ in the same sentence is unaffected: panel A is read straight from
% CANONICAL_NUMBERS.json -> ceilings{}, and raw (0.13178) and rounded-endpoint
% (0.956 - 0.824) both give 0.132 there.
% NOT CHANGED HERE because the correction is a six-site edit and four sites are
% outside this file: 05-Limitations.tex:224 and :225, 03-Apparatus.tex:378, and
% figs/splitref.tex's spread node (:219) and caption (:244). That figure already
% DRAWS train as 0.956 at the resolved line while the prose here still prints
% 0.955, so the disagreement is live. Changing this site alone would replace a
% uniformly wrong number with an inconsistent document. Reported upward instead.
The three splits are not scored on comparable populations either, and the
difference is one of selection. Their raw within-well split-half precision
references, the grey dashes of \cref{fig:splits}a, are $0.956$, $0.824$ and
$0.836$, a spread of $0.132$; \texttt{train}
is filtered at a reliability line the two held-out splits are not, and imposing
that line on all three brings the references to $0.955$, $0.961$ and $0.960$, a
spread of $0.006$. The held-out shortfall therefore cannot be softened by appeal
to a lower noise floor.

% PLACEMENT ONLY. The other retraction in this section is guarded at the
% \needspace above; this one was not, and it broke 13 lines / 2 lines across the
% page. The two lines that landed overleaf are the ones that perform the
% withdrawal, and they arrived with no CORRECTION rule above them, so the reader
% met the retraction as unlabelled prose. \budgettick reserves 5\baselineskip,
% which is not enough: the box measures 247.1pt, 16.6 lines with its padding and
% frame, so 17 is the box entire and it either sets whole here or starts the
% next page whole.
% MEASURED COST, build of 2026-08-12: the guard carries ~2.6 lines of the
% chapter forward. That stops the closing claim of \cref{sec:res-lookup}
% breaking, and starts the coverage definition and the seen-drug-transfer claim
% breaking, all three in 04d-baselines.tex, each of which wants the same guard.
\needspace{17\baselineskip}%

\paragraph{Regression recovers prompt sensitivity} 

Every condition $i$ contributes three scrambled prompts, one per stratum,
each naming a partner drug $p(i,s)$ whose true residual sits at a measured cosine
from the target's. Pooling all three prompts of every condition,
\begin{equation}
 \NIR_{i,s} \;=\; \alpha \;+\; \beta\,
   \cos\!\big(\resid_{p(i,s)},\, \resid_{i}\big) \;+\; \varepsilon_{i,s},
 \qquad s \in \{\texttt{near},\, \texttt{orth},\, \texttt{opposite}\},
\end{equation}
fitted by ordinary least squares with one intercept per condition, so that each
condition is compared only with itself and no part of the slope can come from
conditions differing from one another. The target-drug prompt is never a point on
this line --- only
the three scrambles are --- so nothing in the fit contrasts the model against a
chosen baseline, and $\beta$ is not a gap. It answers a different question: does
the score move with how similar the drug the model was \emph{told} is to the drug
the condition actually received? And $\beta = 0$ is the null by construction, since
a checkpoint that ignores the name emits the same distribution whichever partner
is substituted, putting its three points at one height whatever their cosines
happen to be. No comparator can be blamed for the answer, because none enters it.

Measured, $\beta$ is \ci{+0.361}{+0.249}{+0.472} on \texttt{train},
\ci{+0.321}{+0.207}{+0.434} on \texttt{unseen\_combo} and
\ci{+0.324}{+0.146}{+0.502} on \texttt{unseen\_drug}, clustered on the drug, and
all three still exclude zero clustered on the cell line ($[+0.240, +0.482]$,
$[+0.235, +0.406]$ and $[+0.234, +0.414]$), on $900$, $\num{2685}$ and $597$
points across $300$, $895$ and $199$ conditions. Pooling the conditions instead
of giving each its own intercept moves the estimates barely, to $+0.373$, $+0.336$
and $+0.342$, so unlike the arms of \cref{sec:res-blind} and
\cref{sec:res-epsilon} this one does not depend on the choice. The prediction tracks the drug it is named on every split,
held-out ones included. One restriction attaches to the last: on
\texttt{unseen\_drug} the partners named are overwhelmingly drugs the model
\emph{was} trained on, so the slope there shows it responding to a named trained
drug while scoring a condition whose own drug was withheld, and not that it knows
the withheld drug.


\paragraph{A positive slope beside a chance-level score is not a contradiction}
The two instruments conflict only if they ask the same question, and they do not.
The slope asks whether the prediction moves with the drug it is named. $\NIR$ asks
whether it lands nearer that condition's own truth than the truths of the
$\approx 40$ other drugs in the same cell line. A model passes the first and fails
the second by knowing what a drug does on average and not what it does
\emph{here}, and the size of that shortfall is itself measured:
\cref{sec:res-transfer} puts the share of a drug's residual that transfers across
contexts at \ci{\Tcoef = 0.553}{0.514}{0.593}, leaving
\ci{44.7\%}{40.7\%}{48.6\%} that does not, a remainder mixing cell line, dose,
well and estimator noise. The slope reads the transferable part; winning a ranking
needs the rest. It is also why the lookup wins in \cref{sec:res-lookup}, since
\texttt{drug\_lookup} is that same main effect, estimated straight from data far
more precisely than the model learned it. The condition-level picture agrees: on
trained conditions the cosine between prediction and own truth averages $+0.0389$
(sd $0.1119$, median $+0.0220$), spanning $[-0.3046, +0.3976]$ with $33.7\%$
clearing $|0.1|$ --- predictions pointing the right way on average, and seldom far
enough to win.

\paragraph{One asymmetry attaches to every held-out number in this chapter} The
evaluation's truths are not refitted. The generic is reconstructed from the
build's own fit inventory and verified against the build's \texttt{fit\_digest},
and a mismatch aborts the run. The reliability line selects the \emph{training}
split, so \texttt{train} is by construction the reliability-surviving population,
and the two held-out splits are not filtered at all. In the unseen-combination
arm $\num{499}$ conditions sit below the $+0.1086$ line and $\num{74}$ have a
negative split-half cosine; in the unseen-drug arm, of $\num{199}$ conditions,
$\num{88}$ sit below the line and $\num{18}$ are negative. This is deliberate,
since filtering the evaluation set on its own outcome is the selection this
chapter removed, but it leaves the comparison conservative in the direction that
matters. The held-out arms are scored on their full, noisier populations while
the trained arm is not.

% PLACEMENT ONLY. This box is 23 lines and it was breaking after two of them,
% which is the one split `breakable' is not for: the accent rule, the words "The
% claim." and two lines at the foot of the page, and the whole verdict overleaf.
% It fits on one page with room to spare, so the reservation is the whole box.
\needspace{24\baselineskip}%
\begin{claim}
  Re-encoding the target made the model respond to the combined drug-and-mechanism
  prompt. That much is established and survives every stress: on trained conditions
  it holds clustered on the cell line and on the drug, it holds against chance with
  no comparator at all, and the comparator-free slope excludes zero on all three
  splits under both clusterings --- \ci{+0.361}{+0.249}{+0.472},
  \ci{+0.321}{+0.207}{+0.434} and \ci{+0.324}{+0.146}{+0.502}, fitted with one
  intercept per condition so that each condition is compared only with itself.
  
  Transfer is a different question, and the two instruments answer it with very
  different confidence. The scramble intervention answers it barely. On
  \texttt{unseen\_combo} the gap is $+0.0332$, excluding zero clustered on cell line
  and well, \ci{+0.0332}{+0.0069}{+0.0594}, but not clustered on drug and well,
  \ci{+0.0332}{-0.0031}{+0.0695}, and it is the drug axis a drug-specific claim has
  to survive. The split is weaker than its name --- $30$ of its $34$ drugs and $39$
  of its $42$ cell lines appear in training, and $82\%$ of its conditions recombine
  a seen drug with a seen cell line --- and the effect is a majority tilt rather
  than a uniform one, positive for $20$ of $34$ drugs (sign test $p = 0.39$). The
  model's own $\NIR$ covers chance on both held-out splits. Transfer is
  \notestablished.
  
  The regression answers the same question far more confidently, and the distance
  between the two verdicts is itself the finding. Every point in it is a scrambled
  prompt, so it never asks whether the model is right, only whether it tells wrong
  answers apart --- and by that measure it tells them apart cleanly on every split,
  held-out ones included. Intervening directly on the prompt does not recover an
  effect of the same weight, perhaps because ranking a condition against its neighbours
  needs the part of the residual that does not transfer across contexts, which
  \cref{sec:res-transfer} measures at \ci{44.7\%}{40.7\%}{48.6\%}. 
  
  What the slope on the \emph{unseen-drug} split does and does not add has to be
  stated exactly. At \ci{+0.324}{+0.146}{+0.502} it shows that the prediction still
  tracks the drug \emph{named in the prompt} when the condition's own drug was held
  out. But the drugs named in that regression are the three scrambled partners,
  which on this split are overwhelmingly training drugs, and the real target-drug
  prompt contributes no point to it. The slope is evidence that the model responds
  to named trained drugs in an unseen-drug scoring context, and not evidence that it
  knows the held-out drugs.
  \end{claim}